% !TeX program = pdflatex
% Podręcznik podejmowania decyzji --- wydanie dla inżynierów i programistów --- Basilisk Systems P.S.A.
% Treść jak w basilisk_decyzje.tex, z objaśnieniami w pomarańczowych ramkach „O co chodzi?”.
% Streszcza umowę spółki (psa.tex), umowę akcjonariuszy (basilisk_sha.tex) i dokumenty wewnętrzne.
\documentclass[11pt,a4paper]{report}
\usepackage[T1]{fontenc}
\usepackage[utf8]{inputenc}
\usepackage[english,polish]{babel}
\usepackage{xr-hyper}
\usepackage{basilisk_i18n}
\usepackage{lmodern}
\usepackage{microtype}
\usepackage{setspace}
\usepackage{parskip}
\usepackage{needspace}
\usepackage{tabularx}
\usepackage{float}
\usetikzlibrary{positioning,arrows.meta,shapes.geometric,calc,fit}

\BusinessPlanCompany{Basilisk Systems}
\BusinessPlanWatermarkText{Basilisk Systems}
\BusinessPlanWatermarkColor{black!6}
\BusinessPlanWatermarkFontSize{38}{46}
\BusinessPlanWatermarkAngle{38}
\BusinessPlanAuthorsText{Basilisk Systems --- Założyciele}
\BusinessPlanConfidentialText{Dokument wewnętrzny}
\BusinessPlanFooterText{Basilisk Systems --- podręcznik podejmowania decyzji dla inżynierów}
\BusinessPlanVersionText{Wersja 1.0 --- projekt}

\setstretch{1.1}
\setlist[itemize]{topsep=0.2em,itemsep=0.12em,parsep=0.05em,leftmargin=1.8em}
\setlist[enumerate]{topsep=0.2em,itemsep=0.12em,parsep=0.05em,leftmargin=2.0em}
\setlength{\parindent}{0pt}
\setlength{\parskip}{0.5em}
\setlength{\emergencystretch}{4em}
\renewcommand{\arraystretch}{1.2}
\renewcommand{\thesection}{\arabic{section}}
\captionsetup{font=small,labelfont=bf}

\externaldocument{psa}[psa.pdf]
\externaldocument{basilisk_sha}[basilisk_sha.pdf]

\hypersetup{
  pdftitle={Basilisk Systems P.S.A. --- podręcznik podejmowania decyzji dla inżynierów i programistów},
  pdfauthor={Basilisk Systems}
}

% ---------------------------------------------------------------- style diagramów
\definecolor{cWZ}{HTML}{1F4E79}
\definecolor{cRD}{HTML}{2E7D32}
\definecolor{cDW}{HTML}{8D6E00}
\definecolor{cX}{HTML}{8E2430}
\definecolor{cOC}{HTML}{E65100}
\tikzset{
  >={Stealth[length=2.2mm]},
  every node/.style={font=\small, align=center},
  box/.style={draw, rounded corners=2pt, minimum height=8mm, inner sep=4pt, text width=#1},
  box/.default=3.6cm,
  wz/.style={box=#1, fill=cWZ!12, draw=cWZ},
  wz/.default=3.6cm,
  rd/.style={box=#1, fill=cRD!12, draw=cRD},
  rd/.default=3.6cm,
  dw/.style={box=#1, fill=cDW!14, draw=cDW},
  dw/.default=3.6cm,
  xx/.style={box=#1, fill=cX!10, draw=cX},
  xx/.default=3.6cm,
  q/.style={draw, diamond, aspect=2.4, inner sep=1.5pt, fill=black!4, text width=#1},
  q/.default=3.4cm,
  lbl/.style={font=\scriptsize, fill=white, inner sep=1.5pt},
  line/.style={->, thick},
}
\newtcolorbox{Uwaga}{colback=cX!5, colframe=cX!70, boxrule=0.6pt, arc=2pt, left=6pt, right=6pt, top=4pt, bottom=4pt}
% odesłania do bibliografii (definicje także w basilisk_bib.tex, przez \providecommand)
\makeatletter
\newcommand{\BibItem}[2]{\item[\textbf{[#1]}]\def\@currentlabel{#1}\def\@currentlabelname{#1}\phantomsection\label{bib:#2}}
\makeatother
\newcommand{\Zr}[1]{[\ref{bib:#1}]}
\newtcolorbox{Komentarz}{colback=black!3, colframe=black!45, boxrule=0.5pt, arc=2pt, left=6pt, right=6pt, top=4pt, bottom=4pt, fonttitle=\small\bfseries, title=Komentarz}
\newtcolorbox{OCoChodzi}{colback=cOC!7, colframe=cOC, boxrule=0.8pt, arc=2pt, left=6pt, right=6pt, top=4pt, bottom=4pt, fonttitle=\small\bfseries, coltitle=white, title=O co chodzi?}
\newtcolorbox{Zasada}{colback=cWZ!5, colframe=cWZ!70, boxrule=0.6pt, arc=2pt, left=6pt, right=6pt, top=4pt, bottom=4pt}

\begin{document}

\BusinessPlanTitlePage
  {Basilisk Systems}
  {Wersja 1.0 --- projekt}
  {Podręcznik podejmowania decyzji}
  {Wydanie dla inżynierów i programistów}
  {Kto decyduje, jaką większością i w jakim trybie --- z objaśnieniami „O co chodzi?”}
  {1.0}

\tableofcontents
\newpage

\section{Jak korzystać z podręcznika}

Podręcznik streszcza zasady podejmowania decyzji wynikające z umowy spółki (dalej „US”, wersja 3.4-S), umowy akcjonariuszy (dalej „UA”, wersja 1.8) oraz dokumentów wewnętrznych: polityki bezpieczeństwa informacji, Cennika i wzoru NDA. Odwołania w nawiasach wskazują źródło zasady.

\begin{Uwaga}
Podręcznik nie zmienia żadnej umowy. W razie rozbieżności rozstrzyga treść umowy spółki, a następnie umowy akcjonariuszy. Limity wewnętrzne wiążą dyrektorów wobec Spółki, ale nie wpływają na ważność umów z osobami trzecimi (US §~\ref{art:reprezentacja} ust.~4): dyrektor, który je przekroczy, odpowiada wobec Spółki.
\end{Uwaga}

Materiały źródłowe: umowa spółki \Zr{us}, umowa akcjonariuszy \Zr{ua}, umowy IP \Zr{ip}, Cennik \Zr{cennik}, NDA \Zr{nda}, polityka \Zr{polityka}; podstawy prawne i literatura --- w bibliografii na końcu. Numery paragrafów są klikalne i otwierają odpowiednią umowę (pliki PDF muszą leżeć w tym samym katalogu).

Kolory na diagramach: \textcolor{cWZ}{\textbf{niebieski}} --- akcjonariusze i Walne Zgromadzenie; \textcolor{cRD}{\textbf{zielony}} --- Rada Dyrektorów; \textcolor{cDW}{\textbf{złoty}} --- dyrektor wykonawczy; \textcolor{cX}{\textbf{czerwony}} --- sytuacje szczególne (spór, sąd, konflikt interesów).

\begin{OCoChodzi}
Niniejsze wydanie zawiera pełną treść podręcznika, uzupełnioną o objaśnienia w ramkach „O co chodzi?”. Objaśnienia przedstawiają mechanizmy prawne w kategoriach znanych z inżynierii systemów; nie zmieniają treści umów i nie stanowią porady prawnej.

\textbf{Podstawowe pojęcia:}
\begin{itemize}
  \item \textbf{Umowa spółki (US)} --- dokument nadrzędny, który definiuje organy, ich kompetencje i progi decyzyjne. W razie rozbieżności rozstrzyga jej treść. Zmiana wymaga uchwały z protokołem notarialnym i wpisu do KRS.
  \item \textbf{Umowa akcjonariuszy (UA)} --- umowa między Założycielami; wiąże wyłącznie jej strony.
  \item \textbf{Uchwała} --- formalnie udokumentowana decyzja organu. Czynność wymagająca uchwały, dokonana bez niej, jest wadliwa.
  \item \textbf{Forma dokumentowa} --- dowolny dokument (np. e-mail, PDF), który pozwala ustalić osobę składającą oświadczenie. \textbf{Forma pisemna} --- podpis własnoręczny albo kwalifikowany podpis elektroniczny; zwykły podpis elektroniczny jej nie spełnia.
  \item \textbf{KRS} --- publiczny rejestr przedsiębiorców; część zmian, w szczególności zmiana umowy spółki, wywołuje skutek dopiero z chwilą wpisu.
\end{itemize}
Limity wewnętrzne (budżet, progi kwotowe) mają charakter kontroli organizacyjnej, a nie technicznej blokady: umowa zawarta z ich przekroczeniem jest ważna wobec kontrahenta, a odpowiedzialność ponosi dyrektor, który je naruszył.
\end{OCoChodzi}

% ================================================================
\section{Organy Spółki}

Spółka działa w systemie monistycznym: ma tylko Walne Zgromadzenie i Radę Dyrektorów; nie ma zarządu ani rady nadzorczej (US §~\ref{art:rada} ust.~1).

\begin{figure}[H]
\centering
\begin{tikzpicture}[node distance=9mm and 6mm]
  \node[wz=5.2cm] (akc) {\textbf{Akcjonariusze}\\1 akcja = 1 głos};
  \node[wz=5.2cm, below=of akc] (wz) {\textbf{Walne Zgromadzenie}\\albo uchwała poza WZ (e-mail)};
  \node[rd=5.2cm, below=14mm of wz] (rd) {\textbf{Rada Dyrektorów}\\3--7 dyrektorów, kadencja 3 lata\\Przewodniczący: głos rozstrzygający};
  \node[dw=3.8cm, below=14mm of rd, xshift=-2.6cm] (dw) {\textbf{Dyrektor wykonawczy}\\(co najmniej jeden)\\bieżąca działalność};
  \node[rd=3.8cm, below=14mm of rd, xshift=2.6cm] (dn) {\textbf{Dyrektor niewykonawczy}\\(co najmniej jeden)\\nadzór, umowy z dyr.~wyk.};
  \node[box=2.6cm, right=22mm of rd] (pr) {\textbf{Prokurent}\\(opcjonalnie)};
  \draw[line] (akc) -- (wz);
  \draw[line] (wz) -- node[lbl, right, text width=4.2cm, align=left]{powołuje i odwołuje\\(bezwzględna większość oddanych)} (rd);
  \draw[line] (rd.south) ++(-1.2,0) -- node[lbl, left]{deleguje} (dw.north);
  \draw[line] (rd.south) ++(1.2,0) -- (dn.north);
  \draw[line] (rd) -- node[lbl, above, text width=1.9cm]{zgoda wszystkich dyrektorów} (pr);
  \node[box=10cm, below=8mm of dw.south east, anchor=north, xshift=0.6cm, fill=black!3] (rep) {\textbf{Reprezentacja wobec osób trzecich:} dwóch dyrektorów łącznie albo dyrektor z prokurentem (US §~\ref{art:reprezentacja} ust.~1)};
\end{tikzpicture}
\caption{Struktura organów (US §~\ref{art:rada}--\ref{art:wz}).}
\end{figure}

\begin{itemize}
  \item \textbf{Rada Dyrektorów} prowadzi sprawy Spółki, reprezentuje ją i nadzoruje (US §~\ref{art:rada} ust.~2). Uchwały zapadają bezwzględną większością głosów przy udziale co najmniej połowy dyrektorów; przy równości głosów decyduje Przewodniczący, chyba że sprawa dotyczy jego interesu (US §~\ref{art:rada} ust.~5).
  \item \textbf{Dyrektor wykonawczy} prowadzi bieżącą działalność na podstawie delegacji Rady. Delegacja nie obejmuje spraw wymagających uchwały Rady (US §~\ref{art:rada} ust.~4 i 7).
  \item \textbf{Dyrektor niewykonawczy} --- musi być co najmniej jeden. Reprezentuje Spółkę w umowach i sporach z dyrektorem wykonawczym na podstawie uchwały dyrektorów niewykonawczych (US §~\ref{art:reprezentacja} ust.~3).
  \item \textbf{Regulamin Rady Dyrektorów} (przyjmuje Rada) określa szczegółowy podział kompetencji, limity wydatków i zasady raportowania (US §~\ref{art:rada} ust.~8).
\end{itemize}

% ================================================================
\begin{OCoChodzi}
Struktura organów odpowiada modelowi kontroli dostępu opartej na rolach:
\begin{itemize}
  \item \textbf{Walne Zgromadzenie} --- poziom właścicielski: powołuje i odwołuje dyrektorów oraz zatwierdza decyzje o charakterze fundamentalnym;
  \item \textbf{Rada Dyrektorów} --- organ zarządzający i nadzorczy jednocześnie: prowadzi sprawy Spółki i kontroluje ich wykonanie;
  \item \textbf{dyrektor wykonawczy} --- rola z uprawnieniami delegowanymi przez Radę, ograniczonymi budżetem i regulaminem Rady;
  \item \textbf{dyrektor niewykonawczy} --- niezależna funkcja kontrolna; reprezentuje Spółkę w sprawach z dyrektorem wykonawczym.
\end{itemize}
Reprezentacja łączna (dwóch dyrektorów albo dyrektor z prokurentem) realizuje zasadę czterech oczu: żadna pojedyncza osoba nie zaciąga zobowiązań w imieniu Spółki. Prokura jest szerokim, odwołalnym umocowaniem, udzielanym za zgodą wszystkich dyrektorów.
\end{OCoChodzi}

\begin{Komentarz}
System monistyczny (art.~300$^{73}$ KSH \Zr{ksh}) skraca drogę decyzji: ten sam organ prowadzi sprawy, reprezentuje i nadzoruje. Ceną jest ryzyko, że nadzór stanie się fikcją, jeżeli dyrektor niewykonawczy jest zależny od wykonawczych. W praktyce: delegację dla dyrektora wykonawczego podejmuj uchwałą z datą i zakresem, a dyrektor wykonawczy raportuje Radzie co miesiąc (wykonanie budżetu, umowy powyżej progów, incydenty). Szerzej o modelu PSA: \Zr{opalski}, \Zr{malinowski}, \Zr{ecfr}.
\end{Komentarz}

\section{Kto decyduje? --- drzewo decyzyjne}

Przy każdej nowej sprawie przejdź drzewo od góry. Pierwsza odpowiedź „tak” wskazuje organ i większość.

\begin{figure}[H]
\centering
\begin{tikzpicture}[node distance=7mm and 9mm]
  \node[box=3.2cm, fill=black!8] (start) {\textbf{Nowa sprawa}};
  \node[q, below=of start] (q1) {Pozbawienie prawa poboru?};
  \node[wz=4.4cm, right=of q1] (a1) {\textbf{WZ --- 4/5 głosów}\\US §~\ref{art:emisje} ust.~6};
  \node[q, below=of q1] (q2) {Sprawa z listy\\„75\%” (tabela 1)?};
  \node[wz=4.4cm, right=of q2] (a2) {\textbf{WZ --- 75\% wszystkich głosów w Spółce}\\US §~\ref{art:wz-zastrzezone} ust.~1};
  \node[q, below=of q2] (q3) {Dyrektorzy, sprawozdanie, absolutorium, zysk, wynagrodzenie dyrektorów?};
  \node[wz=4.4cm, right=of q3] (a3) {\textbf{WZ --- bezwzględna większość oddanych}\\US §~\ref{art:wz-zastrzezone} ust.~3};
  \node[q, below=of q3] (q4) {Sprawa z listy Rady (tabela 2)?};
  \node[rd=4.4cm, right=of q4] (a4) {\textbf{Uchwała Rady Dyrektorów}\\US §~\ref{art:rada} ust.~5 i 7};
  \node[dw=4.4cm, below=of q4] (a5) {\textbf{Dyrektor wykonawczy}\\w granicach budżetu i regulaminu Rady};
  \draw[line] (start) -- (q1);
  \foreach \i/\j in {1/2,2/3,3/4} \draw[line] (q\i) -- node[lbl, left]{nie} (q\j);
  \draw[line] (q4) -- node[lbl, left]{nie} (a5);
  \foreach \i in {1,2,3,4} \draw[line] (q\i) -- node[lbl, above]{tak} (a\i);
\end{tikzpicture}
\caption{Drzewo decyzyjne: który organ i jaka większość.}
\end{figure}

\begin{Zasada}
\textbf{„75\% wszystkich głosów w Spółce”} liczy się od wszystkich istniejących głosów, a nie od obecnych. Nieobecny albo wstrzymujący się akcjonariusz działa więc jak głos „przeciw”. \textbf{„Bezwzględna większość oddanych”} oznacza więcej głosów „za” niż łącznie „przeciw” i „wstrzymujących się” spośród oddanych (US §~\ref{art:wz} ust.~7).
\end{Zasada}

\needspace{12\baselineskip}
\subsection*{Tabela 1. Sprawy zastrzeżone dla Walnego Zgromadzenia --- 75\% wszystkich głosów (US §~\ref{art:wz-zastrzezone} ust.~1)}
\begin{tabularx}{\linewidth}{@{}p{0.6cm} X@{}}
\toprule
a) & zmiana umowy spółki (na Walnym Zgromadzeniu, z protokołem notarialnym, i wpis do rejestru) \\
b) & emisja akcji, warrantów, opcji i innych praw do akcji, w tym każda emisja serii F i P \\
c) & nadanie albo cofnięcie statusu Założyciela Funkcjonalnego Serii~F \\
d) & połączenie, podział, przekształcenie, rozwiązanie, likwidacja \\
e) & zbycie albo dzierżawa przedsiębiorstwa lub jego zorganizowanej części; zbycie aktywów ponad 25\% wartości aktywów \\
f) & przeniesienie albo wyłączna, nieodwołalna lub bezterminowa licencja na kluczową własność intelektualną poza zwykłym tokiem działalności \\
g) & nabycie akcji własnych, umorzenie, wypłata z kapitału akcyjnego \\
h) & zobowiązanie, zabezpieczenie albo transakcja poza budżetem powyżej 500\,000 zł \\
i) & zgoda na zbycie akcji w okresie lock-up (12 miesięcy) i zwolnienie akcji z obowiązku zbycia ponad zwolnienie automatyczne \\
j) & kredyt, pożyczka, poręczenie lub podobna umowa z dyrektorem, prokurentem, likwidatorem albo na ich rzecz \\
k) & umowa z akcjonariuszem mającym co najmniej 5\% akcji, dyrektorem, osobą im bliską albo podmiotem powiązanym o wartości ponad 100\,000 zł w roku obrotowym (z wyjątkiem umów Założycieli w zakresie zaangażowania i umów IP z US) \\
\bottomrule
\end{tabularx}

\needspace{12\baselineskip}
\subsection*{Tabela 2. Sprawy wymagające uchwały Rady Dyrektorów (US §~\ref{art:rada} ust.~7)}
\begin{tabularx}{\linewidth}{@{}p{0.6cm} X@{}}
\toprule
1 & przyjęcie i zmiana rocznego budżetu i planu finansowego \\
2 & umowa albo wydatek poza budżetem powyżej 250\,000 zł \\
3 & finansowanie, zabezpieczenie albo rozporządzenie aktywami powyżej 100\,000 zł \\
4 & transakcja z akcjonariuszem, dyrektorem, osobą bliską albo podmiotem powiązanym \\
5 & utworzenie spółki zależnej, działalność w nowej jurysdykcji \\
6 & licencja wyłączna albo przeniesienie istotnej własności intelektualnej (jeżeli nie wymaga WZ --- tabela~1 lit.~f) \\
7 & przyjęcie polityk: bezpieczeństwa informacji, kontroli eksportu, zgodności; Cennika \\
8 & wykonanie uchwał emisyjnych; stwierdzenie wniesienia wkładów \\
9 & zgoda albo odmowa zgody na zbycie lub obciążenie akcji (rozdział~6) \\
10 & stwierdzenie ustania zaangażowania i wskazanie nabywcy akcji niezwolnionych (rozdział~7) \\
11 & pożyczka od Założyciela (UA §~\ref{sha:finansowanie} ust.~2) \\
\bottomrule
\end{tabularx}

% ================================================================
\begin{OCoChodzi}
\textbf{Sposób liczenia głosów.}
\begin{itemize}
  \item Próg „75\% wszystkich głosów” odnosi się do wszystkich akcji w Spółce, a nie do akcji reprezentowanych przy głosowaniu. Brak głosu ma więc w praktyce taki sam skutek jak głos przeciw.
  \item Przy podziale 50/20/10/10/10 Założyciel~A może zablokować każdą uchwałę 75\% (bez niego można zebrać najwyżej 50\%), ale nie przeprowadzi jej sam. Założyciele A i~B łącznie mają 70\%, więc potrzebny jest jeszcze trzeci Założyciel.
  \item Bezwzględna większość głosów oddanych oznacza, że głosów „za” jest więcej niż łącznie głosów „przeciw” i „wstrzymujących się”.
\end{itemize}
\textbf{Skutki braku zgody zależą od jej źródła.} Brak uchwały wymaganej ustawą powoduje nieważność czynności. Brak zgody wymaganej wyłącznie umową spółki nie wpływa na ważność czynności wobec kontrahenta, ale rodzi odpowiedzialność dyrektora wobec Spółki.
\end{OCoChodzi}

\begin{Komentarz}
Najczęstszy błąd to mylenie „75\% wszystkich głosów w Spółce” z 75\% obecnych. Przy podziale 50/20/10/10/10 Założyciel~A sam nie przeprowadzi uchwały 75\%, ale może każdą zablokować; A z~B mają 70\% --- potrzebny jest jeszcze trzeci Założyciel. Skutki pominięcia zgody są różne: czynność dokonana bez uchwały wymaganej \emph{ustawą} jest nieważna (art.~17 §~1 KSH \Zr{ksh}), a bez zgody wymaganej \emph{tylko umową spółki} --- ważna wobec kontrahenta, lecz dyrektor odpowiada wobec Spółki (US §~\ref{art:reprezentacja} ust.~4).
\end{Komentarz}

\section{Progi kwotowe}

Ta sama czynność może wymagać różnych decyzji zależnie od kwoty. Kwoty dotyczą jednej czynności albo kilku czynności powiązanych.

\begin{figure}[H]
\centering
\begin{tikzpicture}[x=0.0165cm, y=1cm]
  % oś
  \draw[->] (0,0) -- (640,0) node[right, font=\scriptsize]{tys. zł};
  \foreach \v in {0,100,250,500} {\draw (\v,0.08) -- (\v,-0.08) node[below, font=\scriptsize]{\v};}
  % pas 1: wydatek poza budżetem
  \node[anchor=east, font=\scriptsize, text width=3.4cm, align=right] at (-8,1.2) {umowa / wydatek\\poza budżetem};
  \fill[cDW!25] (0,0.85) rectangle (250,1.55);  \node[font=\scriptsize] at (125,1.2) {dyrektor wykonawczy};
  \fill[cRD!25] (250,0.85) rectangle (500,1.55); \node[font=\scriptsize] at (375,1.2) {Rada Dyrektorów};
  \fill[cWZ!25] (500,0.85) rectangle (630,1.55); \node[font=\scriptsize] at (565,1.2) {WZ 75\%};
  % pas 2: finansowanie / zabezpieczenie / aktywa
  \node[anchor=east, font=\scriptsize, text width=3.4cm, align=right] at (-8,2.3) {finansowanie, zabezpie-\\czenie, zbycie aktywów\\(WZ: tylko poza budżetem)};
  \fill[cDW!25] (0,1.95) rectangle (100,2.65);  \node[font=\scriptsize] at (50,2.3) {dyr.~wyk.};
  \fill[cRD!25] (100,1.95) rectangle (500,2.65); \node[font=\scriptsize] at (300,2.3) {Rada Dyrektorów};
  \fill[cWZ!25] (500,1.95) rectangle (630,2.65); \node[font=\scriptsize] at (565,2.3) {WZ 75\%};
  % pas 3: w budżecie
  \node[anchor=east, font=\scriptsize, text width=3.4cm, align=right] at (-8,3.4) {wydatek w ramach\\przyjętego budżetu};
  \fill[cDW!25] (0,3.05) rectangle (630,3.75); \node[font=\scriptsize] at (315,3.4) {dyrektor wykonawczy (w granicach regulaminu Rady)};
  \foreach \v in {100,250,500} \draw[densely dotted, black!50] (\v,0) -- (\v,3.85);
\end{tikzpicture}
\caption{Kto zatwierdza czynność w zależności od kwoty (US §~\ref{art:rada} ust.~7, §~\ref{art:wz-zastrzezone} ust.~1 lit.~h).}
\end{figure}

Niezależnie od kwoty: transakcja z akcjonariuszem, dyrektorem albo podmiotem powiązanym wymaga uchwały Rady (dyrektor mający konflikt interesów nie głosuje, ale liczy się do kworum; gdy uchwały nie można podjąć, rozstrzyga WZ), a umowa kredytu lub pożyczki z dyrektorem --- także zgody WZ 75\% (tabela~1 lit.~j). Umowa z akcjonariuszem mającym co najmniej 5\%, dyrektorem albo podmiotem powiązanym ponad 100\,000 zł rocznie wymaga WZ 75\% (lit.~k). Zbycie aktywów ponad 25\% ich wartości zawsze wymaga WZ 75\% (lit.~e).

% ================================================================
\begin{OCoChodzi}
Progi kwotowe przypisują kompetencję do organu według wartości czynności. Przykłady:
\begin{itemize}
  \item zakup serwera obliczeniowego za 180\,000 zł ujęty w budżecie --- dyrektor wykonawczy;
  \item ten sam zakup poza budżetem za 300\,000 zł --- uchwała Rady Dyrektorów (próg 250\,000 zł);
  \item leasing sprzętu za 150\,000 zł --- finansowanie, a więc uchwała Rady (próg 100\,000 zł);
  \item licencja na oprogramowanie symulacyjne za 600\,000 zł poza budżetem --- Walne Zgromadzenie, 75\% wszystkich głosów;
  \item umowa z podmiotem, w którym dyrektor ma udziały --- zawsze uchwała Rady, a powyżej 100\,000 zł rocznie także Walne Zgromadzenie.
\end{itemize}
Podział jednej czynności na kilka mniejszych w celu zejścia poniżej progu jest obejściem umowy spółki; czynności powiązane ocenia się łącznie. Najskuteczniejszym sposobem ograniczenia liczby uchwał jest szczegółowy budżet, który z góry zatwierdza planowane wydatki.
\end{OCoChodzi}

\begin{Komentarz}
Progi dotyczą także czynności powiązanych: podzielenie jednej umowy na części, aby zejść poniżej progu, jest obejściem i nie chroni dyrektora przed odpowiedzialnością. Im dokładniej Rada przyjmie budżet, tym mniej uchwał ad hoc --- budżet jest głównym narzędziem przekazania decyzji dyrektorowi wykonawczemu. Progi nie podlegają waloryzacji; warto przeglądać je co 2--3 lata. \textbf{Zasada legalna:} progi 100/250/500 tys.\ zł są zapisane w umowie spółki i ich złagodzenie (podwyższenie) wymaga zmiany umowy --- 75\% wszystkich głosów, notariusz, KRS. Akcjonariusze mogą natomiast zwykłą uchwałą \emph{zaostrzyć} limity w stosunkach wewnętrznych (np. obniżyć próg dla danego rodzaju umów albo wymagać zgody Rady na dodatkowe czynności), a Rada może ustalić limity dla dyrektora wykonawczego w regulaminie Rady (US §~\ref{art:rada} ust.~8). Takie ograniczenia wiążą dyrektorów wobec Spółki, ale nie wpływają na ważność umów z osobami trzecimi (US §~\ref{art:reprezentacja} ust.~4); art.~300$^{53}$, art.~300$^{77}$ §~2 i art.~17 §~3 KSH \Zr{ksh}, \Zr{malinowski}.
\end{Komentarz}

\section{Jak podjąć uchwałę}

\subsection*{Uchwała akcjonariuszy poza Walnym Zgromadzeniem (e-mail)}

Domyślny tryb dla większości spraw. Wyjątek: uchwały wymagające protokołu notarialnego (w szczególności zmiana umowy spółki) podejmuje się na Walnym Zgromadzeniu (US §~\ref{art:wz} ust.~1).

\begin{figure}[H]
\centering
\begin{tikzpicture}[node distance=6mm and 5mm]
  \node[wz=2.9cm] (s1) {\textbf{1. Projekt}\\Rada, Przewodniczący albo akcjonariusz z $\ge$5\% głosów rozsyła projekt z uzasadnieniem};
  \node[wz=2.9cm, right=of s1] (s2) {\textbf{2. Głosowanie}\\e-mailem z adresu z rejestru; termin: 3 dni robocze -- 21 dni};
  \node[wz=2.9cm, right=of s2] (s3) {\textbf{3. Podjęcie}\\z upływem terminu albo wcześniej, gdy jest większość wszystkich głosów};
  \node[wz=2.9cm, right=of s3] (s4) {\textbf{4. Wynik}\\Przewodniczący stwierdza wynik i zawiadamia e-mailem; wpis do księgi protokołów};
  \draw[line] (s1) -- (s2); \draw[line] (s2) -- (s3); \draw[line] (s3) -- (s4);
\end{tikzpicture}
\caption{Uchwała poza Walnym Zgromadzeniem (US §~\ref{art:wz} ust.~2--3 i 8).}
\end{figure}

\begin{itemize}
  \item Krótszy termin niż 3 dni robocze jest możliwy za zgodą wszystkich akcjonariuszy. Głos oddany po terminie się nie liczy.
  \item Głosowanie \emph{na piśmie} (papier) wymaga zgody wszystkich akcjonariuszy na ten tryb; głosowanie e-mailem takiej zgody nie wymaga.
  \item Walne Zgromadzenie bez formalnego zwołania jest możliwe, jeżeli reprezentowane są wszystkie akcje i nikt nie zgłasza sprzeciwu (US §~\ref{art:wz} ust.~4). Udział zdalny reguluje regulamin WZ.
\end{itemize}

\subsection*{Uchwała Rady Dyrektorów}

\begin{enumerate}
  \item Przewodniczący (albo dyrektor) rozsyła projekt; głosowanie może odbyć się na posiedzeniu, na piśmie albo e-mailem, jeżeli da się zidentyfikować uczestników i utrwalić głosy (US §~\ref{art:rada} ust.~6).
  \item \textbf{Konflikt interesów:} dyrektor, którego interes jest sprzeczny z interesem Spółki, ujawnia to i nie głosuje, ale liczy się go do kworum, jeżeli uczestniczy w posiedzeniu albo głosowaniu; gdy wskutek wyłączeń uchwały nie można podjąć, rozstrzyga WZ (US §~\ref{art:rada} ust.~6).
  \item Kworum: co najmniej połowa powołanych dyrektorów. Większość: bezwzględna. Przy równości --- głos Przewodniczącego.
  \item Uchwałę zapisuje się i przechowuje w dokumentacji Spółki.
\end{enumerate}

% ================================================================
\begin{OCoChodzi}
Uchwała podejmowana drogą elektroniczną to sformalizowana procedura głosowania asynchronicznego:
\begin{itemize}
  \item projekt z uzasadnieniem rozsyła uprawniony podmiot, a głos oddaje się z adresu e-mail wpisanego do rejestru akcjonariuszy, co zapewnia identyfikację głosującego;
  \item uchwała zostaje podjęta z upływem terminu albo wcześniej --- z chwilą oddania wymaganej większości wszystkich głosów;
  \item głos oddany po terminie nie jest uwzględniany;
  \item treść uchwały i wynik głosowania wpisuje się do księgi protokołów, która pełni funkcję trwałego, niemodyfikowalnego rejestru decyzji.
\end{itemize}
W Radzie Dyrektorów dyrektor pozostający w konflikcie interesów nie uczestniczy w głosowaniu, ale jest uwzględniany przy ustalaniu kworum. Jeżeli wskutek wyłączeń uchwały nie można podjąć, sprawę rozstrzyga Walne Zgromadzenie.
\end{OCoChodzi}

\begin{Komentarz}
Wiadomość e-mail wysłana z adresu wpisanego do rejestru akcjonariuszy spełnia wymóg formy dokumentowej (art.~77$^{2}$ KC \Zr{kc}). Projekt, głosy i zawiadomienie o wyniku trzymaj w jednym wątku, a uchwałę wpisz do księgi protokołów (art.~300$^{100}$ KSH). Zawiadomienie o wyniku wysyłaj od razu --- od niego zależy pewność, kiedy uchwała została podjęta, co ma znaczenie m.in. dla terminów jej zaskarżenia.
\end{Komentarz}

\section{Zgoda na zbycie akcji}

Każde zbycie albo obciążenie akcji wymaga zgody Spółki, którą wydaje Rada Dyrektorów (US §~\ref{art:zbywanie}). Spółka \textbf{musi} wyrazić zgodę, jeżeli nie zachodzi żadna z pięciu przyczyn odmowy.

\begin{figure}[H]
\centering
\begin{tikzpicture}[node distance=7mm and 8mm]
  \node[box=4.2cm, fill=black!8] (z) {\textbf{Zgłoszenie do Spółki}\\nabywca, cena, istotne warunki};
  \node[q=3.3cm, below=of z] (lock) {Lock-up 12 mies.\\od wpisu / Dnia Przyznania F?};
  \node[wz=3.4cm, right=of lock] (lockwz) {zgoda WZ 75\%\\(albo czekać)};
  \node[rd=4.2cm, below=of lock] (rd) {\textbf{Rada Dyrektorów}\\14 dni na zgodę albo odmowę\\(forma dokumentowa, e-mail)};
  \node[rd=3.4cm, below=10mm of rd] (no) {\textbf{Odmowa}\\tylko z przyczyn a)--e), z ich wskazaniem};
  \node[rd=2.8cm, left=6mm of no] (ok) {\textbf{Zgoda}\\$\to$ wpis nabywcy do rejestru};
  \node[xx=3.2cm, right=6mm of no] (cisza) {\textbf{Milczenie}\\$\ne$ zgoda};
  \node[rd=3.4cm, below=of no] (e) {Przyczyna e) (konkurent): Spółka w 30 dni wskazuje innego nabywcę na tych samych warunkach --- inaczej musi udzielić zgody};
  \node[xx=3.2cm, below=of cisza] (w7) {wezwanie Rady\\+ 7 dni};
  \node[wz=3.2cm, below=of w7] (ua) {uchwała akcjonariuszy poza WZ --- 14 dni (tylko z przyczyn a--e)};
  \node[xx=3.2cm, below=of ua] (sad) {sąd: żądanie złożenia oświadczenia o zgodzie (art.~64 KC)};
  \draw[line] (z) -- (lock);
  \draw[line] (lock) -- node[lbl, above]{tak} (lockwz);
  \draw[line] (lockwz) |- node[lbl, pos=0.3, right]{po zgodzie WZ} (rd);
  \draw[line] (lock) -- node[lbl, left]{nie} (rd);
  \draw[line] (rd) -- (ok); \draw[line] (rd) -- (no); \draw[line] (rd) -- (cisza);
  \draw[line] (no) -- (e);
  \draw[line] (cisza) -- (w7); \draw[line] (w7) -- (ua); \draw[line] (ua) -- node[lbl, right]{brak zgody\\i brak odmowy} (sad);
\end{tikzpicture}
\caption{Procedura zgody na zbycie akcji (US §~\ref{art:zbywanie} ust.~1--4).}
\end{figure}

\textbf{Przyczyny odmowy (zamknięty katalog, US §~\ref{art:zbywanie} ust.~2):} a)~sankcje albo sprzeczność z przepisami o kontroli inwestycji, eksportu lub bezpieczeństwie państwa; b)~brak wymaganej zgody organu (można dać zgodę warunkową); c)~lock-up bez zgody WZ; d)~akcje objęte obowiązkiem zbycia, a nabywca nie przejął obowiązków (pełnomocnictwo i oferta); e)~nabywca jest konkurentem (z wyjątkiem inwestycji portfelowych do 5\%).

Zgody nie wymaga nabycie akcji w wykonaniu obowiązku zbycia ani zbycie w egzekucji (US §~\ref{art:zbywanie} ust.~5).

% ================================================================
\begin{OCoChodzi}
Zgoda na zbycie akcji jest decyzją związaną, a nie uznaniową: Rada Dyrektorów musi jej udzielić, chyba że zachodzi jedna z pięciu przyczyn odmowy wskazanych w umowie spółki. Katalog jest zamknięty; odmowa z innej przyczyny jest nieskuteczna.
\begin{itemize}
  \item \textbf{Lock-up} --- przez 12 miesięcy zbycie wymaga zgody Walnego Zgromadzenia (75\%).
  \item \textbf{Brak odpowiedzi nie oznacza zgody.} Umowa przewiduje procedurę eskalacji: wezwanie Rady z dodatkowym terminem 7 dni, następnie uchwała akcjonariuszy, a w ostateczności --- wyrok sądu, który zastępuje oświadczenie Spółki.
  \item \textbf{Weryfikacja nabywcy} pod kątem sankcji i kontroli inwestycji jest obowiązkowym elementem procedury; w spółce z sektora obronnego niewłaściwy nabywca może zagrozić koncesji.
\end{itemize}
\end{OCoChodzi}

\begin{Komentarz}
Zgoda na zbycie jest decyzją związaną: Rada nie ocenia ceny ani osoby nabywcy poza katalogiem z US §~\ref{art:zbywanie} ust.~2, a odmowa bez wskazania podstawy nie jest skuteczną odmową. Przyczyny z lit.~a--b wymagają sprawdzenia nabywcy na listach sankcyjnych \Zr{sankcje}, \Zr{usank} i pod kątem kontroli inwestycji \Zr{uki}, \Zr{fdi} --- wynik sprawdzenia warto opisać w uchwale. Jeżeli Spółka bezpodstawnie nie udziela zgody, prawomocny wyrok zastępuje jej oświadczenie (art.~64 KC \Zr{kc}, art.~1047 KPC \Zr{kpc}).
\end{Komentarz}

\section{Vesting i ustanie zaangażowania}

\subsection*{Harmonogram zwalniania akcji (US §~\ref{art:obowiazek-zbycia} ust.~2)}

\begin{figure}[H]
\centering
\begin{tikzpicture}[x=0.245cm, y=0.042cm]
  \draw[->] (0,0) -- (51,0) node[right, font=\scriptsize]{miesiące};
  \draw[->] (0,0) -- (0,108) node[above, font=\scriptsize]{\% akcji zwolnionych};
  \foreach \m in {0,6,12,24,36,48} \draw (\m,1.5) -- (\m,-1.5) node[below, font=\scriptsize]{\m};
  \foreach \p in {25,50,75,100} {\draw (0.3,\p) -- (-0.3,\p) node[left, font=\scriptsize]{\p}; \draw[black!12] (0,\p) -- (50,\p);}
  % Założyciel Pierwotny
  \draw[very thick, cWZ] (0,25) -- (6,25) -- (6,37.5) -- (36,100) -- (49,100);
  \fill[cWZ] (0,25) circle (0.25 and 1.6);
  \node[font=\scriptsize, cWZ, anchor=west, text width=4.2cm, align=left] at (1,90) {\textbf{Założyciel Pierwotny:} 25\% od razu, cliff 6 mies., potem co miesiąc do 36 mies.};
  % Seria F
  \draw[very thick, cRD, dashed] (0,0) -- (12,0) -- (12,25) -- (48,100) -- (49,100);
  \node[font=\scriptsize, cRD, anchor=west, text width=4.2cm, align=left] at (26,40) {\textbf{Seria F:} cliff 12 mies. (25\%), potem co miesiąc do 48 mies.};
\end{tikzpicture}
\caption{Zwalnianie akcji z obowiązku zbycia. Do chwili zwolnienia akcje dają pełne prawa (US §~\ref{art:obowiazek-zbycia} ust.~3).}
\end{figure}

\subsection*{Co się dzieje przy odejściu (US §~\ref{art:obowiazek-zbycia} ust.~4--5, UA §~\ref{sha:ustanie}--\ref{sha:wykonanie})}

\begin{figure}[H]
\centering
\begin{tikzpicture}[node distance=7mm and 6mm]
  \node[box=4.6cm, fill=black!8] (ev) {\textbf{Ustanie zaangażowania}\\(rozwiązanie umowy, rezygnacja, śmierć, niezdolność)};
  \node[rd=4.6cm, below=of ev] (st) {\textbf{Rada stwierdza ustanie} w 30 dni (bez udziału zainteresowanego): dzień i liczba akcji niezwolnionych};
  \node[q=3.1cm, below=of st] (typ) {Rodzaj odejścia?};
  \node[wz=3.6cm, below=10mm of typ] (u) {\textbf{Odejście Uzasadnione} (bez istotnej przyczyny, w tym niepowołanie; uzgodnione)\\zwalnia się dodatkowo 12 mies.};
  \node[wz=3.4cm, left=6mm of u] (d) {\textbf{Śmierć / trwała niezdolność}\\zwalniają się \textbf{wszystkie} akcje};
  \node[xx=3.4cm, right=6mm of u] (z) {\textbf{Odejście Zwykłe}\\bez dodatkowego zwolnienia};
  \node[rd=6.2cm, below=12mm of u] (n) {Akcje niezwolnione $\to$ zbycie po \textbf{cenie emisyjnej} (akcje za aport: wartość wkładu, max wartość godziwa): najpierw pozostali Założyciele Pierwotni pro rata w 30 dni, potem nabywca wskazany przez Radę w 2 mies.; po 6 mies. obowiązek wygasa (bieg zawieszony w sporze)};
  \node[wz=6.2cm, below=of n] (wz75) {WZ może (75\%) zwolnić także pozostałe akcje};
  \draw[line] (ev) -- (st); \draw[line] (st) -- (typ);
  \draw[line] (typ) -| (d); \draw[line] (typ) -- (u); \draw[line] (typ) -| (z);
  \draw[line] (u) -- (n); \draw[line] (z) |- (n);
  \draw[line] (n) -- (wz75);
\end{tikzpicture}
\caption{Ustanie zaangażowania i los akcji niezwolnionych.}
\end{figure}

\begin{itemize}
  \item Zwolnienie przy śmierci, niezdolności i Odejściu Uzasadnionym następuje \textbf{automatycznie} z mocy umowy spółki --- nie wymaga głosowania (UA §~\ref{sha:ustanie} ust.~4).
  \item Umowę zbycia w imieniu odchodzącego może zawrzeć pełnomocnik (Spółka albo dyrektor niewykonawczy) na podstawie nieodwołalnego pełnomocnictwa; nabywca płaci w 14 dni, wpis do rejestru dopiero po zapłacie (UA §~\ref{sha:wykonanie}).
  \item Po Odejściu Zwykłym obowiązuje zakaz konkurencji przez 12 miesięcy; zakaz aktywnego nakłaniania pracowników i klientów obowiązuje 12 miesięcy przy każdym odejściu (UA §~\ref{sha:konkurencja}).
\end{itemize}

% ================================================================
\begin{OCoChodzi}
Mechanizm zwalniania akcji różni się od typowego programu opcji. Założyciel od początku jest właścicielem wszystkich swoich akcji i wykonuje z nich pełnię praw (głos, dywidenda). Harmonogram określa jedynie, jaka część akcji podlega obowiązkowi odsprzedaży w razie wcześniejszego zakończenia współpracy.

Przykład orientacyjny dla Założyciela Pierwotnego (według wykresu): 25\% akcji zwolnionych od dnia zawiązania, 37,5\% po 6 miesiącach, następnie liniowo do 100\% po 36 miesiącach (ok.~2,08 punktu procentowego miesięcznie).
\begin{itemize}
  \item Odejście Zwykłe w 18. miesiącu: zwolnione ok.~62,5\% akcji; pozostałe ok.~37,5\% podlega odsprzedaży.
  \item Odejście Uzasadnione w tym samym momencie: dodatkowe 12 miesięcy, czyli zwolnione ok.~87,5\%.
  \item Śmierć albo trwała niezdolność do pracy: zwolnienie wszystkich akcji.
\end{itemize}
Cena odkupu odpowiada cenie emisyjnej, a dla akcji objętych za wkład niepieniężny --- wartości wkładu, nie wyższej niż wartość godziwa. Pracowników i współpracowników obejmuje odrębny program motywacyjny (seria~P): opcje nabywane przez 4 lata, z rocznym okresem karencji.
\end{OCoChodzi}

\begin{Komentarz}
Odkup akcji niezwolnionych po cenie emisyjnej to rozwiązanie rynkowe; dla akcji objętych za aport IP umowa przewiduje cenę równą wartości wkładu (nie wyższą niż wartość godziwa), co ogranicza ryzyko zarzutu, że Spółka zatrzymuje wkład bez ekwiwalentu. Mimo to sąd może badać jego zgodność z zasadami współżycia społecznego (art.~58 i 353$^{1}$ KC \Zr{kc}); dlatego umowa uzasadnia cenę mechanizmem przyjętym przez wszystkich Założycieli przy zawiązaniu Spółki (UA §~\ref{sha:wykonanie} ust.~2). Przy zbyciu poniżej wartości rynkowej trzeba ocenić skutki w podatku dochodowym \Zr{pit}. Rada powinna stwierdzić ustanie zaangażowania w 30 dni --- spóźnienie nie znosi obowiązku zbycia, ale osłabia pozycję Spółki w ewentualnym sporze.
\end{Komentarz}

\section{Umowy i spory z dyrektorami}

\begin{figure}[H]
\centering
\begin{tikzpicture}[node distance=7mm and 8mm]
  \node[box=4.4cm, fill=black!8] (u) {\textbf{Umowa albo spór Spółki z dyrektorem}};
  \node[q=3.0cm, below=of u] (q) {Dyrektor wykonawczy?};
  \node[rd=4.6cm, below left=9mm and -10mm of q] (dn) {Dyrektor niewykonawczy na podstawie uchwały dyrektorów niewykonawczych\\\textbf{albo} pełnomocnik z uchwały akcjonariuszy};
  \node[wz=4.6cm, below right=9mm and -10mm of q] (pe) {\textbf{Pełnomocnik} powołany uchwałą akcjonariuszy do zawarcia, zmiany, wykonania i rozwiązania umowy (może być poza WZ; nie strona umowy; zalecenie: nie dyrektor)};
  \node[xx=9.8cm, anchor=north] (poz) at ($(q |- pe.south)+(0,-5mm)$) {\textbf{Pożyczka, kredyt, poręczenie z dyrektorem:} uchwała Rady + zgoda WZ 75\% + pełnomocnik z uchwały akcjonariuszy (US §~\ref{art:wz-zastrzezone} ust.~1 lit.~j, UA §~\ref{sha:finansowanie} ust.~2)};
  \draw[line] (u) -- (q);
  \draw[line] (q) -- node[lbl, left]{tak} (dn);
  \draw[line] (q) -- node[lbl, right]{nie} (pe);
\end{tikzpicture}
\caption{Reprezentacja Spółki w sprawach z dyrektorem (US §~\ref{art:reprezentacja} ust.~3).}
\end{figure}

Dotyczy to w szczególności umów przeniesienia praw IP (Umowa nr 1 i nr 2) z Założycielami będącymi dyrektorami, umów wykonawczych do obowiązku zbycia i potwierdzeń kwartalnych. Jedna uchwała akcjonariuszy może powołać pełnomocnika do wszystkich umów danego rodzaju --- do ich zawarcia, zmiany, wykonania i rozwiązania oraz innych czynności z nimi związanych (US §~\ref{art:reprezentacja} ust.~3). Ten szeroki zakres jest potrzebny, bo szczególny tryb reprezentacji obejmuje każdą czynność związaną z umową, np. aneks do umowy IP albo jej rozwiązanie przy odejściu Założyciela (\Zr{sn4218}).

% ================================================================
\begin{OCoChodzi}
Obowiązuje zasada rozdziału ról: dyrektor nie może reprezentować Spółki w umowie, której jest stroną, ani w umowie z innym dyrektorem. Spółkę reprezentuje wówczas osoba niezależna --- pełnomocnik powołany uchwałą akcjonariuszy albo dyrektor niewykonawczy działający na podstawie uchwały dyrektorów niewykonawczych.

\textbf{Znaczenie dla zespołu technicznego:} umowy przeniesienia praw do kodu i dokumentacji (Umowa nr~1 i nr~2) zawierane z Założycielami pełniącymi funkcję dyrektorów muszą być podpisane w imieniu Spółki przez takiego pełnomocnika. Uchybienie grozi nieważnością umowy, a w konsekwencji --- brakiem praw Spółki do kluczowej technologii. Jest to jeden z pierwszych punktów weryfikowanych w badaniu due diligence.
\end{OCoChodzi}

\begin{Komentarz}
Na tle art.~210 KSH przyjmuje się, że umowa zawarta z członkiem organu z naruszeniem szczególnych zasad reprezentacji jest nieważna, a szczególny tryb obejmuje także zmianę, wypowiedzenie i rozwiązanie takiej umowy (\Zr{sn4218}, \Zr{sa48518}, \Zr{jara}); to samo ryzyko dotyczy art.~300$^{79}$ KSH. Najbezpieczniej podjąć przy zawiązaniu jedną uchwałę akcjonariuszy powołującą pełnomocnika do umów IP i umów wykonawczych z dyrektorami (US §~\ref{art:final} ust.~6) i pilnować, by pełnomocnik nie był stroną danej umowy ani innym dyrektorem --- według orzecznictwa do art.~210 KSH członek organu nie może reprezentować spółki jako pełnomocnik w umowie z innym członkiem tego organu. Wyjątkiem jest droga z art.~300$^{79}$ §~3 KSH: dyrektor niewykonawczy działający na podstawie uchwały dyrektorów niewykonawczych (US §~\ref{art:reprezentacja} ust.~3).
\end{Komentarz}

\section{Zmiana umowy spółki}

Zmiana umowy spółki wymaga uchwały akcjonariuszy podjętej na Walnym Zgromadzeniu z protokołem notarialnym i wpisu do rejestru przedsiębiorców.

\begin{figure}[H]
\centering
\begin{tikzpicture}[node distance=6mm and 7mm]
  \node[box=4.6cm, fill=black!8] (p) {\textbf{Projekt zmiany umowy\\{\normalfont Rada albo akcjonariusz}}};
  \node[wz=4.6cm, below=of p] (wz) {\textbf{Walne Zgromadzenie\\{\normalfont z notariuszem (nie e-mailem)}}};
  \node[wz=4.6cm, below=of wz] (g) {\textbf{75\% wszystkich głosów w Spółce\\{\normalfont US §~\ref{art:wz-zastrzezone} ust.~1 lit.~a}}};
  \node[xx=3.6cm, left=8mm of g] (pp) {\textbf{Wyłączenie prawa poboru?}\\$\to$ 4/5 głosów\\(US §~\ref{art:emisje} ust.~6)};
  \node[xx=3.6cm, right=8mm of g] (ob) {\textbf{Więcej obowiązków akcjonariusza?} (lock-up, vesting)\\$\to$ + zgoda każdego, kogo dotyczy (US §~\ref{art:final} ust.~1)};
  \node[wz=4.6cm, below=of g] (n) {\textbf{Protokół notarialny}};
  \node[rd=4.6cm, below=of n] (k) {\textbf{Wpis do KRS\\{\normalfont dopiero wtedy zmiana działa}}};
  \draw[line] (p) -- (wz); \draw[line] (wz) -- (g); \draw[line] (g) -- (n); \draw[line] (n) -- (k);
  \draw[line, dashed] (g) -- (pp); \draw[line, dashed] (g) -- (ob);
\end{tikzpicture}
\caption{Zmiana umowy spółki (US §~\ref{art:wz} ust.~1, §~\ref{art:wz-zastrzezone} ust.~1 lit.~a, §~\ref{art:final} ust.~1).}
\end{figure}

\begin{OCoChodzi}
Zmiana umowy spółki przebiega w trzech etapach, z których każdy jest konieczny:
\begin{itemize}
  \item uchwała Walnego Zgromadzenia podjęta większością 75\% wszystkich głosów;
  \item protokół notarialny, który nadaje uchwale wymaganą formę;
  \item wpis do KRS --- dopiero z tą chwilą zmiana wywołuje skutek prawny.
\end{itemize}
Zmiana nakładająca na akcjonariuszy nowe obowiązki (np. wydłużenie okresu lock-up) wymaga ponadto zgody każdego akcjonariusza, którego dotyczy; sama większość nie wystarcza. Przypomina to zmianę niekompatybilną wstecz, która wymaga akceptacji każdego, na kogo wpływa.
\end{OCoChodzi}

\begin{Komentarz}
Przy podziale 50/20/10/10/10: Założyciel~A (50\%) sam nie zmieni umowy, ale może zablokować każdą zmianę; A z~B (70\%) to za mało --- potrzebny jest jeszcze co najmniej jeden Założyciel z 10\%. Bez A zmiany nie da się przeprowadzić (pozostali mają razem 50\%).
\end{Komentarz}

\section{Emisje akcji}

\begin{figure}[H]
\centering
\begin{tikzpicture}[node distance=7mm and 6mm]
  \node[q=3.4cm] (q) {Seria A, F albo P,\\w limitach?};
  \node[wz=4.0cm, right=10mm of q] (zm) {\textbf{Najpierw zmiana umowy} (WZ 75\%, notariusz, KRS) --- nowa seria albo wyższy limit};
  \node[wz=4.0cm, below=of q] (em) {\textbf{Uchwała emisyjna WZ 75\%}\\maksymalna liczba akcji, cena minimalna, termin};
  \node[wz=3.4cm, left=10mm of em] (pp) {Wyłączenie prawa poboru?\\$\to$ \textbf{4/5 głosów}};
  \node[rd=4.0cm, below=of em] (rd) {\textbf{Rada wykonuje:} oferty, umowy objęcia, przyjęcie wkładów, rejestr, zgłoszenia};
  \draw[line] (q) -- node[lbl, above]{nie} (zm);
  \draw[line] (q) -- node[lbl, left]{tak} (em);
  \draw[line] (zm) |- (em);
  \draw[line, dashed] (em) -- (pp);
  \draw[line] (em) -- (rd);
\end{tikzpicture}
\caption{Emisja akcji (US §~\ref{art:emisje}--\ref{art:seria-p}).}
\end{figure}

\begin{itemize}
  \item \textbf{Limity:} łącznie 2\,000\,000 akcji; seria F do 250\,000 (i do 20\% akcji); seria P do 176\,471; termin emisji do 31 grudnia 2036 r. (US §~\ref{art:emisje} ust.~2--3).
  \item \textbf{Seria F --- dwie uchwały WZ 75\%:} najpierw nadanie statusu Założyciela Funkcjonalnego (rola, zaangażowanie, maksymalna liczba akcji, Dzień Przyznania), potem emisja (US §~\ref{art:seria-f}).
  \item \textbf{Seria P:} regulamin programu zatwierdza WZ; Rada prowadzi program i rekomenduje uczestników, ale nie emituje akcji samodzielnie (US §~\ref{art:seria-p}).
  \item Cena emisyjna nie niższa niż 0,01 zł; Rada nie może jej obniżyć poniżej minimum z uchwały (US §~\ref{art:emisje} ust.~5 i 7).
\end{itemize}

% ================================================================
\begin{OCoChodzi}
Kluczowe pojęcia związane z emisją akcji:
\begin{itemize}
  \item \textbf{Limity} --- umowa spółki określa łączny limit 2\,000\,000 akcji i limity poszczególnych serii; ich przekroczenie wymaga zmiany umowy spółki.
  \item \textbf{Rozwodnienie} --- po emisji udział procentowy dotychczasowych akcjonariuszy maleje; wartość ich akcji może jednak wzrosnąć, jeżeli inwestor obejmuje akcje po wyższej wycenie.
  \item \textbf{Prawo poboru} --- pierwszeństwo objęcia nowych akcji w proporcji do posiadanych, pozwalające utrzymać udział. Umowa spółki wyłącza je dla serii~P (program motywacyjny) i~S (zamiana pożyczek i umów SAFE).
  \item \textbf{Podział kompetencji} --- o emisji decyduje Walne Zgromadzenie; Rada Dyrektorów wyłącznie ją wykonuje, w granicach uchwały.
\end{itemize}
\end{OCoChodzi}

\begin{Komentarz}
Typowa kolejność rundy inwestorskiej: term sheet $\to$ zmiana umowy spółki tworząca serię inwestorską (Walne Zgromadzenie, notariusz) $\to$ uchwała emisyjna $\to$ umowy objęcia akcji $\to$ wpis do rejestru akcjonariuszy $\to$ zgłoszenie do rejestru przedsiębiorców; równolegle zmiana umowy akcjonariuszy (UA §~\ref{sha:finansowanie} ust.~3). Przy wkładzie w postaci pracy (seria~F) pamiętaj o terminie wniesienia wkładu --- w ciągu 3 lat (art.~300$^{9}$ §~1 w zw. z art.~300$^{104}$ §~2 KSH).
\end{Komentarz}

\section{Dokumenty wewnętrzne --- kto przyjmuje}

\begin{tabularx}{\linewidth}{@{}>{\raggedright\arraybackslash}p{4.0cm} >{\raggedright\arraybackslash}p{3.3cm} X@{}}
\toprule
\textbf{Dokument} & \textbf{Przyjmuje} & \textbf{Uwagi} \\
\midrule
Roczny budżet i plan finansowy & Rada Dyrektorów & wyznacza granice działania dyrektora wykonawczego \\
Regulamin Rady Dyrektorów & Rada Dyrektorów & podział kompetencji, limity, raportowanie, konflikt interesów \\
Regulamin Walnego Zgromadzenia & akcjonariusze & udział zdalny; pierwszy przyjmują Założyciele \\
Polityka bezpieczeństwa informacji & Rada Dyrektorów & wyznacza Odpowiedzialnego; każdy Użytkownik potwierdza zapoznanie się \\
Cennik wynagrodzeń za rozwiązania projektowe & Rada Dyrektorów & waloryzacja GUS co roku 1 marca; zmiana na niekorzyść nie działa wstecz \\
Regulamin programu motywacyjnego (seria P) & Walne Zgromadzenie & warunki, ograniczenia zbywania, skutki ustania współpracy \\
Wzór NDA & Rada Dyrektorów & zawiera się \textbf{przed} ujawnieniem informacji (UA §~\ref{sha:poufnosc} ust.~2 lit.~b); podpisuje dwóch dyrektorów albo dyrektor z prokurentem \\
Lista osób z dostępem do informacji Kontrolowanych & Rada Dyrektorów & polityka §~3; przegląd dostępów co kwartał \\
\bottomrule
\end{tabularx}

% ================================================================
\begin{OCoChodzi}
Dla zespołu technicznego najważniejszym dokumentem wewnętrznym jest polityka bezpieczeństwa informacji. Jej kluczowe wymagania:
\begin{itemize}
  \item konta Spółki z uwierzytelnianiem dwuskładnikowym, hasła w menedżerze haseł, klucze i tokeny w menedżerze sekretów --- nigdy w kodzie ani w wiadomościach;
  \item kod przechowywany wyłącznie w repozytoriach organizacji Spółki, bez kopii na kontach prywatnych;
  \item zakaz wprowadzania informacji Poufnych i Kontrolowanych do ogólnodostępnych narzędzi, w tym narzędzi sztucznej inteligencji, poza narzędziami zatwierdzonymi;
  \item przetwarzanie informacji Kontrolowanych wyłącznie na urządzeniach Spółki;
  \item rejestrowanie dostępu do repozytoriów i systemów oraz możliwość przeglądu tych rejestrów.
\end{itemize}
\textbf{Kontrola eksportu.} Udostępnienie kodu, modelu albo danych osobie przebywającej za granicą --- również przez repozytorium, aktualizację zdalną albo dostęp zdalny --- może stanowić wywóz technologii wymagający zezwolenia. Przed takim udostępnieniem należy uzyskać ocenę Odpowiedzialnego.

Wymagania te mają również znaczenie prawne: kod i know-how korzystają z ochrony jako tajemnica przedsiębiorstwa tylko wtedy, gdy Spółka podejmuje realne działania w celu zachowania ich poufności.
\end{OCoChodzi}

\begin{Komentarz}
Polityka bezpieczeństwa, rejestr dostępów i zawarte NDA są dowodem „niezbędnych działań” w rozumieniu art.~11 ust.~2 UZNK \Zr{uznk}; bez nich ochrona tajemnicy przedsiębiorstwa może upaść (\Zr{sn30400}). Kary umowne w NDA i umowach z twórcami należą się bez wykazywania szkody, ale podlegają miarkowaniu (\Zr{sn6103}); wobec pracowników na etacie kary umownej nie stosuje się (\Zr{sn12007}). Kwoty kar i Cennika są waloryzowane wskaźnikiem GUS \Zr{guscpi}.
\end{Komentarz}

\section{Impas i spory}

\begin{figure}[H]
\centering
\begin{tikzpicture}[node distance=6mm and 6mm]
  \node[box=3.0cm, fill=black!8] (s) {\textbf{Spór między akcjonariuszami}};
  \node[box=3.0cm, right=of s] (n) {negocjacje\\30 dni od wezwania};
  \node[xx=3.0cm, right=of n] (m) {mediacja\\(ośrodek z UA)};
  \node[xx=3.0cm, right=of m] (sd) {sąd powszechny\\siedziby Spółki};
  \draw[line] (s) -- (n); \draw[line] (n) -- (m); \draw[line] (m) -- (sd);
  \node[wz=7.2cm, below=9mm of n.south east, anchor=north] (imp) {\textbf{Impas:} WZ dwa razy (w odstępie $\ge$30 dni) nie podejmuje uchwały o emisji (lit.~b), transakcji ponad 500\,000 zł (lit.~h) albo w sprawie z §~13 ust.~3 --- mediacja; w sprawach z §~13 ust.~3 po 60 dniach obowiązuje status quo (UA)};
  \draw[line] (imp.north) ++(1.7,0) -- (m.south);
\end{tikzpicture}
\caption{Rozwiązywanie sporów (UA §~\ref{sha:spory}). Zabezpieczenie sądowe jest dostępne zawsze.}
\end{figure}

% ================================================================
\begin{OCoChodzi}
Umowa akcjonariuszy przewiduje stopniową procedurę rozwiązywania sporów: negocjacje (30 dni), mediacja, a dopiero następnie sąd.

\textbf{Status quo po 60 dniach mediacji} zapewnia ciągłość działania Spółki w razie impasu: dotychczasowi dyrektorzy zostają ponownie powołani, sprawozdanie finansowe zostaje zatwierdzone, zysk pozostaje w Spółce, a wynagrodzenia nie ulegają zmianie. Spółka funkcjonuje więc w bezpiecznej konfiguracji do czasu porozumienia.

\textbf{Zabezpieczenie sądowe} jest dostępne niezależnie od mediacji. W razie ujawnienia kodu albo naruszenia zakazu konkurencji Spółka może niezwłocznie wystąpić do sądu o zakazanie dalszych naruszeń.
\end{OCoChodzi}

\begin{Komentarz}
Ugoda zawarta przed mediatorem i zatwierdzona przez sąd ma moc ugody sądowej \Zr{kpc}; mediacja nie wstrzymuje biegu terminów do zaskarżenia uchwał. Zabezpieczenie roszczeń (art.~730 i nast. KPC) jest dostępne niezależnie od mediacji --- ma znaczenie zwłaszcza przy naruszeniu poufności lub zakazu konkurencji.
\end{Komentarz}

\section{Kalendarz terminów}

\begin{longtable}{@{}>{\raggedright\arraybackslash}p{6.2cm} >{\raggedright\arraybackslash}p{4.4cm} >{\raggedright\arraybackslash}p{3.6cm}@{}}
\toprule
\textbf{Czynność} & \textbf{Termin} & \textbf{Źródło} \\
\midrule
\endhead
Uchwała akcjonariuszy poza WZ & 3 dni robocze -- 21 dni na głos & US §~\ref{art:wz} ust.~3 \\
Zgoda albo odmowa zgody na zbycie akcji & 14 dni od zgłoszenia & US §~\ref{art:zbywanie} ust.~1 \\
Dodatkowy termin po wezwaniu Rady & 7 dni & US §~\ref{art:zbywanie} ust.~1 \\
Uchwała akcjonariuszy w sprawie zgody & 14 dni od rozesłania projektu & US §~\ref{art:zbywanie} ust.~1 \\
Wskazanie nabywcy po odmowie (konkurent) & 30 dni od odmowy & US §~\ref{art:zbywanie} ust.~3 \\
Lock-up Założycieli Pierwotnych / serii F & 12 mies. od wpisu / Dnia Przyznania & US §~\ref{art:zbywanie} ust.~4 \\
Stwierdzenie ustania zaangażowania & 30 dni od dowiedzenia się & UA §~\ref{sha:ustanie} ust.~2 \\
Pierwszeństwo nabycia przez pozostałych Założycieli & 30 dni od zawiadomienia przez Radę & US §~\ref{art:obowiazek-zbycia} ust.~4 \\
Wskazanie nabywcy pozostałych akcji przez Radę & kolejne 2 mies. & US §~\ref{art:obowiazek-zbycia} ust.~4 \\
Wygaśnięcie obowiązku zbycia & 6 mies. od ustania (bieg zawieszony w sporze) & US §~\ref{art:obowiazek-zbycia} ust.~4 \\
Status quo po bezskutecznej mediacji (sprawy z US §~\ref{art:wz-zastrzezone} ust.~3) & 60 dni od wszczęcia mediacji & UA §~\ref{sha:spory} ust.~3 \\
Zapłata ceny akcji niezwolnionych & 14 dni od umowy zbycia & UA §~\ref{sha:wykonanie} ust.~4 \\
Wydanie utworów po odejściu & 14 dni & Umowa nr~2 §~7 ust.~4, polityka §~8 \\
Odebranie dostępów & najpóźniej ostatni dzień; przy sporze od razu & polityka §~8 \\
Zgłoszenie incydentu / naruszenia danych & 24 godziny / 72 godziny do UODO & polityka §~7 \\
Waloryzacja kar i Cennika & co roku 1 marca & UA §~\ref{sha:kary} ust.~2, Cennik §~5 \\
Negocjacje przed mediacją & 30 dni od wezwania & UA §~\ref{sha:spory} ust.~1 \\
Kadencja Rady Dyrektorów & 3 lata obrotowe & US §~\ref{art:rada} ust.~3 \\
Termin emisji serii F i P & do 31 grudnia 2036 r. & US §~\ref{art:emisje} ust.~3 \\
\bottomrule
\end{longtable}

\begin{OCoChodzi}
Terminy należy traktować jak kamienie milowe w harmonogramie projektu: każdy powinien mieć właściciela i zostać wpisany do kalendarza albo systemu zgłoszeń z przypomnieniem. Zasady obliczania (art.~111--112 KC): dnia zdarzenia nie wlicza się (termin biegnący od wezwania doręczonego 3 marca liczy się od 4 marca); „dni” oznaczają dni kalendarzowe, chyba że wskazano „dni robocze”; termin oznaczony w miesiącach kończy się z upływem dnia, który datą odpowiada początkowemu dniowi terminu. Przekroczenie terminu przez Radę co do zasady nie oznacza zgody, lecz otwiera procedurę eskalacji.
\end{OCoChodzi}

% Wspólna bibliografia Basilisk Systems --- % Wspólna bibliografia Basilisk Systems --- % Wspólna bibliografia Basilisk Systems --- \input{basilisk_bib} w podręczniku decyzji i komentarzu do umowy spółki.
% Odesłanie w tekście: \Zr{klucz} -> [ETYKIETA] z linkiem do pozycji.
\makeatletter
\providecommand{\BibItem}[2]{\item[\textbf{[#1]}]\def\@currentlabel{#1}\def\@currentlabelname{#1}\phantomsection\label{bib:#2}}
\makeatother
\providecommand{\Zr}[1]{[\ref{bib:#1}]}
\newenvironment{BibList}{\begin{list}{}{\setlength{\labelwidth}{2.6cm}\setlength{\leftmargin}{2.9cm}\setlength{\labelsep}{0.3cm}\setlength{\itemsep}{0.15em}\setlength{\parsep}{0pt}\small}}{\end{list}}

\chapter*{Bibliografia}
\addcontentsline{toc}{chapter}{Bibliografia}

Akty prawne podano według miejsca pierwotnej publikacji oraz ostatniego tekstu jednolitego (stan: listopad 2026), z dopiskiem „ze zm.”; przed użyciem należy sprawdzić, czy nie ogłoszono nowszego tekstu w Internetowym Systemie Aktów Prawnych (isap.sejm.gov.pl) albo w Dzienniku Ustaw (eli.gov.pl). Tezy orzeczeń podano w skrócie; przed powołaniem się na nie w sporze należy sprawdzić pełny tekst.

\section*{A. Akty prawne krajowe}
\begin{BibList}
\BibItem{KSH}{ksh} Ustawa z dnia 15 września 2000 r. --- Kodeks spółek handlowych (Dz.U. 2000 nr 94 poz. 1037 ze zm.; t.j. Dz.U. 2024 poz. 18 ze zm., w tym ustawą z dnia 23 stycznia 2026 r., Dz.U. 2026 poz. 176 --- zmiany rejestru akcjonariuszy od 18 lutego 2027 r.); przepisy o prostej spółce akcyjnej: dział IA, art.~300$^{1}$--300$^{134}$.
\BibItem{Nowela PSA}{psa2019} Ustawa z dnia 19 lipca 2019 r. o zmianie ustawy --- Kodeks spółek handlowych oraz niektórych innych ustaw (Dz.U. 2019 poz. 1655) --- wprowadzenie prostej spółki akcyjnej; ustawa w pierwotnym brzmieniu wchodziła w życie 1 marca 2020 r., a przepisy o PSA stosuje się od 1 lipca 2021 r. (po przesunięciu terminu).
\BibItem{KC}{kc} Ustawa z dnia 23 kwietnia 1964 r. --- Kodeks cywilny (Dz.U. 1964 nr 16 poz. 93; t.j. Dz.U. 2026 poz. 795 ze zm.).
\BibItem{KRO}{kro} Ustawa z dnia 25 lutego 1964 r. --- Kodeks rodzinny i opiekuńczy (Dz.U. 1964 nr 9 poz. 59; t.j. Dz.U. 2026 poz. 236 ze zm.).
\BibItem{KPC}{kpc} Ustawa z dnia 17 listopada 1964 r. --- Kodeks postępowania cywilnego (Dz.U. 1964 nr 43 poz. 296; t.j. Dz.U. 2026 poz. 468 ze zm.).
\BibItem{KP}{kp} Ustawa z dnia 26 czerwca 1974 r. --- Kodeks pracy (Dz.U. 1974 nr 24 poz. 141; t.j. Dz.U. 2025 poz. 277 ze zm.).
\BibItem{PrAut}{praut} Ustawa z dnia 4 lutego 1994 r. o prawie autorskim i prawach pokrewnych (Dz.U. 1994 nr 24 poz. 83; t.j. Dz.U. 2025 poz. 24 ze zm.).
\BibItem{PWP}{pwp} Ustawa z dnia 30 czerwca 2000 r. --- Prawo własności przemysłowej (Dz.U. 2001 nr 49 poz. 508; t.j. Dz.U. 2023 poz. 1170 ze zm.).
\BibItem{UZNK}{uznk} Ustawa z dnia 16 kwietnia 1993 r. o zwalczaniu nieuczciwej konkurencji (Dz.U. 1993 nr 47 poz. 211; t.j. Dz.U. 2026 poz. 85 ze zm.).
\BibItem{PIT}{pit} Ustawa z dnia 26 lipca 1991 r. o podatku dochodowym od osób fizycznych (Dz.U. 1991 nr 80 poz. 350; t.j. Dz.U. 2026 poz. 592 ze zm.).
\BibItem{CIT}{cit} Ustawa z dnia 15 lutego 1992 r. o podatku dochodowym od osób prawnych (Dz.U. 1992 nr 21 poz. 86 ze zm.).
\BibItem{UoR}{uor} Ustawa z dnia 29 września 1994 r. o rachunkowości (Dz.U. 1994 nr 121 poz. 591; t.j. Dz.U. 2026 poz. 522 ze zm.).
\BibItem{UKRS}{ukrs} Ustawa z dnia 20 sierpnia 1997 r. o Krajowym Rejestrze Sądowym (Dz.U. 1997 nr 121 poz. 769; t.j. Dz.U. 2025 poz. 869 ze zm.).
\BibItem{UONI}{uoni} Ustawa z dnia 5 sierpnia 2010 r. o ochronie informacji niejawnych (Dz.U. 2010 nr 182 poz. 1228; t.j. Dz.U. 2025 poz. 1209 ze zm.).
\BibItem{UOTS}{uots} Ustawa z dnia 29 listopada 2000 r. o obrocie z zagranicą towarami, technologiami i usługami o znaczeniu strategicznym dla bezpieczeństwa państwa, a także dla utrzymania międzynarodowego pokoju i bezpieczeństwa (Dz.U. 2000 nr 119 poz. 1250; t.j. Dz.U. 2023 poz. 1582 ze zm., w tym Dz.U. 2026 poz. 471 --- definicje i zezwolenia, m.in. pomoc techniczna).
\BibItem{UWT}{uwt} Ustawa z dnia 13 czerwca 2019 r. o wykonywaniu działalności gospodarczej w zakresie wytwarzania i obrotu materiałami wybuchowymi, bronią, amunicją oraz wyrobami i technologią o przeznaczeniu wojskowym lub policyjnym (Dz.U. 2019 poz. 1214; t.j. Dz.U. 2023 poz. 1743 ze zm.); koncesja na obrót technologią o przeznaczeniu wojskowym (art.~7), wymogi wobec wspólników z co najmniej 20\% akcji (art.~10 i 17)..
\BibItem{UKI}{uki} Ustawa z dnia 24 lipca 2015 r. o kontroli niektórych inwestycji (Dz.U. 2015 poz. 1272; t.j. Dz.U. 2026 poz. 47 ze zm.).
\BibItem{USank}{usank} Ustawa z dnia 13 kwietnia 2022 r. o szczególnych rozwiązaniach w zakresie przeciwdziałania wspieraniu agresji na Ukrainę oraz służących ochronie bezpieczeństwa narodowego (Dz.U. 2022 poz. 835; t.j. Dz.U. 2025 poz. 514 ze zm.).
\BibItem{USyg}{usyg} Ustawa z dnia 14 czerwca 2024 r. o ochronie sygnalistów (Dz.U. 2024 poz. 928).
\BibItem{UODO}{uodo} Ustawa z dnia 10 maja 2018 r. o ochronie danych osobowych (Dz.U. 2018 poz. 1000 ze zm.).
\BibItem{PrLot}{prlot} Ustawa z dnia 3 lipca 2002 r. --- Prawo lotnicze (Dz.U. 2002 nr 130 poz. 1112; t.j. Dz.U. 2025 poz. 1431 ze zm.).
\BibItem{PKD 2025}{pkd} Rozporządzenie Rady Ministrów z dnia 18 grudnia 2024 r. w sprawie Polskiej Klasyfikacji Działalności (PKD) (Dz.U. 2024 poz. 1936).
\BibItem{WT}{wt} Rozporządzenie Rady Ministrów z dnia 17 września 2019 r. w sprawie klasyfikacji rodzajów materiałów wybuchowych, broni, amunicji oraz wyrobów i technologii o przeznaczeniu wojskowym lub policyjnym, na których wytwarzanie lub obrót jest wymagane uzyskanie koncesji (Dz.U. 2019 poz. 1888) --- obejmuje m.in. bezzałogowe statki powietrzne i ich systemy (WT~V pkt~3).
\end{BibList}

\section*{B. Akty prawne Unii Europejskiej}
\begin{BibList}
\BibItem{RODO}{rodo} Rozporządzenie Parlamentu Europejskiego i Rady (UE) 2016/679 z dnia 27 kwietnia 2016 r. w sprawie ochrony osób fizycznych w związku z przetwarzaniem danych osobowych i w sprawie swobodnego przepływu takich danych oraz uchylenia dyrektywy 95/46/WE (ogólne rozporządzenie o ochronie danych) (Dz.Urz. UE L 119 z 4.5.2016, s.~1).
\BibItem{Dual-use}{dualuse} Rozporządzenie Parlamentu Europejskiego i Rady (UE) 2021/821 z dnia 20 maja 2021 r. ustanawiające unijny system kontroli wywozu, pośrednictwa, pomocy technicznej, tranzytu i transferu produktów podwójnego zastosowania (wersja przekształcona) (Dz.Urz. UE L 206 z 11.6.2021, s.~1, ze zm.).
\BibItem{EDF}{edf} Rozporządzenie Parlamentu Europejskiego i Rady (UE) 2021/697 z dnia 29 kwietnia 2021 r. ustanawiające Europejski Fundusz Obronny i uchylające rozporządzenie (UE) 2018/1092 (Dz.Urz. UE L 170 z 12.5.2021, s.~149).
\BibItem{UAS 945}{uas945} Rozporządzenie delegowane Komisji (UE) 2019/945 z dnia 12 marca 2019 r. w sprawie bezzałogowych systemów powietrznych oraz operatorów bezzałogowych systemów powietrznych z państw trzecich (Dz.Urz. UE L 152 z 11.6.2019, s.~1).
\BibItem{UAS 947}{uas947} Rozporządzenie wykonawcze Komisji (UE) 2019/947 z dnia 24 maja 2019 r. w sprawie przepisów i procedur dotyczących eksploatacji bezzałogowych statków powietrznych (Dz.Urz. UE L 152 z 11.6.2019, s.~45).
\BibItem{FDI}{fdi} Rozporządzenie Parlamentu Europejskiego i Rady (UE) 2019/452 z dnia 19 marca 2019 r. ustanawiające ramy monitorowania bezpośrednich inwestycji zagranicznych w Unii --- uchylone z dniem 17 stycznia 2028 r. przez \Zr{fdi2026}.
\BibItem{FDI 2026}{fdi2026} Rozporządzenie Parlamentu Europejskiego i Rady (UE) 2026/1386 z dnia 17 czerwca 2026 r. w sprawie monitorowania inwestycji zagranicznych w Unii i uchylające rozporządzenie (UE) 2019/452 (Dz.Urz. UE L 2026/1386 z 26.6.2026); stosuje się od 17 stycznia 2028 r. (niektóre przepisy od 16 lipca 2026 r.); obejmuje m.in. spółki rozwijające produkty podwójnego zastosowania, wojskowe i drony.
\BibItem{TTBER}{ttber} Rozporządzenie Komisji (UE) 2026/877 z dnia 16 kwietnia 2026 r. w sprawie stosowania art.~101 ust.~3 TFUE do kategorii porozumień o transferze technologii (Dz.Urz. UE L 2026/877 z 21.4.2026); obowiązuje od 1 maja 2026 r. do 30 kwietnia 2038 r.; zastąpiło rozporządzenie (UE) nr 316/2014 --- licencje na technologię, know-how i oprogramowanie.
\BibItem{AI Act}{aiact} Rozporządzenie Parlamentu Europejskiego i Rady (UE) 2024/1689 z dnia 13 czerwca 2024 r. w sprawie ustanowienia zharmonizowanych przepisów dotyczących sztucznej inteligencji (akt w sprawie sztucznej inteligencji) (Dz.Urz. UE L 2024/1689 z 12.7.2024); systemy wyłącznie wojskowe są wyłączone (art.~2 ust.~3), systemy podwójnego zastosowania --- objęte.
\BibItem{Sankcje}{sankcje} Rozporządzenie Rady (UE) nr 833/2014 z dnia 31 lipca 2014 r. dotyczące środków ograniczających w związku z działaniami Rosji destabilizującymi sytuację na Ukrainie (Dz.Urz. UE L 229 z 31.7.2014, s.~1) oraz rozporządzenie Rady (UE) nr 269/2014 z dnia 17 marca 2014 r. (Dz.Urz. UE L 78 z 17.3.2014, s.~6) (środki ograniczające).
\BibItem{TSD}{tsd} Dyrektywa Parlamentu Europejskiego i Rady (UE) 2016/943 z dnia 8 czerwca 2016 r. w sprawie ochrony niejawnego know-how i niejawnych informacji handlowych (tajemnic przedsiębiorstwa) przed ich bezprawnym pozyskiwaniem, wykorzystywaniem i ujawnianiem (Dz.Urz. UE L 157 z 15.6.2016, s.~1).
\BibItem{WBD}{wbd} Dyrektywa Parlamentu Europejskiego i Rady (UE) 2019/1937 z dnia 23 października 2019 r. w sprawie ochrony osób zgłaszających naruszenia prawa Unii.
\end{BibList}

\section*{C. Orzecznictwo}
\begin{BibList}
\BibItem{SN 61/03}{sn6103} Uchwała składu siedmiu sędziów SN z dnia 6 listopada 2003 r., III CZP 61/03, OSNC 2004, nr 5, poz. 69 (zasada prawna) --- kara umowna należy się także wtedy, gdy wierzyciel nie poniósł szkody; dłużnik może żądać jej miarkowania (art.~484 §~2 KC).
\BibItem{SN 304/00}{sn30400} Wyrok SN z dnia 3 października 2000 r., I CKN 304/00, OSNC 2001, nr 4, poz. 59 --- wykorzystanie przez pracownika informacji, co do których przedsiębiorca nie podjął niezbędnych działań w celu zachowania ich poufności, należy traktować jako wykorzystanie powszechnej wiedzy, do której przedsiębiorca nie ma ustawowych uprawnień.
\BibItem{SN 120/07}{sn12007} Wyrok SN z dnia 8 stycznia 2008 r., II PK 120/07 --- odpowiedzialność pracownika za szkodę regulują wyczerpująco przepisy Kodeksu pracy; kara umowna (art.~483 KC) zastrzeżona w stosunku pracy nie może być dochodzona; dopuszczalna jest natomiast w umowie o zakazie konkurencji po ustaniu stosunku pracy (SN powołuje wyrok z 10 października 2003 r., I PK 528/02, OSNP 2004, nr 19, poz. 336).
\BibItem{SN 42/18}{sn4218} Uchwała SN z dnia 7 września 2018 r., III CZP 42/18, OSNC 2019, nr 6, poz. 64 --- art.~210 §~1 KSH stosuje się do zgody spółki z o.o. na zmianę umowy spółki komandytowej, której wspólnikiem jest także członek zarządu; przepis obejmuje wszystkie czynności związane z umową między spółką a członkiem zarządu (zawarcie, zmianę, wypowiedzenie, rozwiązanie, potrącenie, ugodę), a do jego zastosowania wystarcza potencjalna kolizja interesów. Przydatna przy wykładni art.~300$^{79}$ KSH.
\BibItem{SA 485/18}{sa48518} Wyrok Sądu Apelacyjnego w Warszawie z dnia 9 kwietnia 2019 r., V ACa 485/18 --- członek zarządu nie może jako pełnomocnik reprezentować spółki z o.o. w umowie z innym członkiem zarządu (art.~210 §~1 KSH; sąd powołuje doktrynę oraz uchwały SN III CZP 71/18 i III CZP 42/18); skutkiem naruszenia jest bezwzględna nieważność czynności (art.~58 §~1 KC).
\end{BibList}

\section*{D. Literatura}
\begin{BibList}
\BibItem{Malinowski}{malinowski} P. Malinowski (red.), \emph{Prosta spółka akcyjna. Komentarz}, Warszawa 2021.
\BibItem{Opalski}{opalski} A. Opalski, \emph{Prosta spółka akcyjna --- nowy typ spółki handlowej}, cz.~I: „Przegląd Prawa Handlowego” 2019, nr 11, s.~5--15; cz.~II: tamże, nr 12, s.~4--16.
\BibItem{ECFR}{ecfr} \emph{The Simple Joint Stock Company: Emergence of a New Close Company in Poland}, „European Company and Financial Law Review” 2024, t.~21, DOI 10.1515/ecfr-2024-0021.
\BibItem{Jara}{jara} Z. Jara (red.), \emph{Kodeks spółek handlowych. Komentarz}, C.H. Beck (wydanie aktualne).
\BibItem{Gniewek}{gniewek} E. Gniewek, P. Machnikowski (red.), \emph{Kodeks cywilny. Komentarz}, C.H. Beck (wydanie aktualne).
\BibItem{Barta}{barta} J. Barta, R. Markiewicz (red.), \emph{Ustawa o prawie autorskim i prawach pokrewnych. Komentarz}, Wolters Kluwer (wydanie aktualne).
\BibItem{Kostański}{kostanski} P. Kostański (red.), \emph{Prawo własności przemysłowej. Komentarz}, C.H. Beck (wydanie aktualne).
\BibItem{Szwaja}{szwaja} J. Szwaja (red.), \emph{Ustawa o zwalczaniu nieuczciwej konkurencji. Komentarz}, C.H. Beck (wydanie aktualne).
\end{BibList}

\section*{E. Materiały urzędowe}
\begin{BibList}
\BibItem{GUS PKD}{guspkd} Główny Urząd Statystyczny, \emph{Polska Klasyfikacja Działalności PKD 2025} (klasyfikacje.stat.gov.pl).
\BibItem{GUS CPI}{guscpi} Komunikaty Prezesa GUS w sprawie średniorocznego wskaźnika cen towarów i usług konsumpcyjnych ogółem, ogłaszane w Dzienniku Urzędowym RP „Monitor Polski” (podstawa waloryzacji kar i Cennika).
\end{BibList}

\section*{F. Dokumenty Spółki}
\begin{BibList}
\BibItem{US}{us} Umowa spółki Basilisk Systems P.S.A., wersja 3.4-S (\texttt{psa.pdf}; wersja angielska \texttt{psa\_en.pdf}).
\BibItem{UA}{ua} Umowa akcjonariuszy, wersja 1.8 (\texttt{basilisk\_sha.pdf}).
\BibItem{IP}{ip} Umowy przeniesienia praw autorskich (Umowa nr~1 i nr~2), wersja 1.6 (\texttt{basilisk\_ip.pdf}).
\BibItem{Cennik}{cennik} Cennik wynagrodzeń za rozwiązania projektowe (\texttt{basilisk\_cennik.pdf}).
\BibItem{NDA}{nda} Wzór umowy o zachowaniu poufności (\texttt{basilisk\_nda.pdf}).
\BibItem{Polityka}{polityka} Polityka bezpieczeństwa informacji (\texttt{basilisk\_polityka.pdf}).
\BibItem{Podręcznik}{podrecznik} Podręcznik podejmowania decyzji (\texttt{basilisk\_decyzje.pdf}).
\end{BibList}
 w podręczniku decyzji i komentarzu do umowy spółki.
% Odesłanie w tekście: \Zr{klucz} -> [ETYKIETA] z linkiem do pozycji.
\makeatletter
\providecommand{\BibItem}[2]{\item[\textbf{[#1]}]\def\@currentlabel{#1}\def\@currentlabelname{#1}\phantomsection\label{bib:#2}}
\makeatother
\providecommand{\Zr}[1]{[\ref{bib:#1}]}
\newenvironment{BibList}{\begin{list}{}{\setlength{\labelwidth}{2.6cm}\setlength{\leftmargin}{2.9cm}\setlength{\labelsep}{0.3cm}\setlength{\itemsep}{0.15em}\setlength{\parsep}{0pt}\small}}{\end{list}}

\chapter*{Bibliografia}
\addcontentsline{toc}{chapter}{Bibliografia}

Akty prawne podano według miejsca pierwotnej publikacji oraz ostatniego tekstu jednolitego (stan: listopad 2026), z dopiskiem „ze zm.”; przed użyciem należy sprawdzić, czy nie ogłoszono nowszego tekstu w Internetowym Systemie Aktów Prawnych (isap.sejm.gov.pl) albo w Dzienniku Ustaw (eli.gov.pl). Tezy orzeczeń podano w skrócie; przed powołaniem się na nie w sporze należy sprawdzić pełny tekst.

\section*{A. Akty prawne krajowe}
\begin{BibList}
\BibItem{KSH}{ksh} Ustawa z dnia 15 września 2000 r. --- Kodeks spółek handlowych (Dz.U. 2000 nr 94 poz. 1037 ze zm.; t.j. Dz.U. 2024 poz. 18 ze zm., w tym ustawą z dnia 23 stycznia 2026 r., Dz.U. 2026 poz. 176 --- zmiany rejestru akcjonariuszy od 18 lutego 2027 r.); przepisy o prostej spółce akcyjnej: dział IA, art.~300$^{1}$--300$^{134}$.
\BibItem{Nowela PSA}{psa2019} Ustawa z dnia 19 lipca 2019 r. o zmianie ustawy --- Kodeks spółek handlowych oraz niektórych innych ustaw (Dz.U. 2019 poz. 1655) --- wprowadzenie prostej spółki akcyjnej; ustawa w pierwotnym brzmieniu wchodziła w życie 1 marca 2020 r., a przepisy o PSA stosuje się od 1 lipca 2021 r. (po przesunięciu terminu).
\BibItem{KC}{kc} Ustawa z dnia 23 kwietnia 1964 r. --- Kodeks cywilny (Dz.U. 1964 nr 16 poz. 93; t.j. Dz.U. 2026 poz. 795 ze zm.).
\BibItem{KRO}{kro} Ustawa z dnia 25 lutego 1964 r. --- Kodeks rodzinny i opiekuńczy (Dz.U. 1964 nr 9 poz. 59; t.j. Dz.U. 2026 poz. 236 ze zm.).
\BibItem{KPC}{kpc} Ustawa z dnia 17 listopada 1964 r. --- Kodeks postępowania cywilnego (Dz.U. 1964 nr 43 poz. 296; t.j. Dz.U. 2026 poz. 468 ze zm.).
\BibItem{KP}{kp} Ustawa z dnia 26 czerwca 1974 r. --- Kodeks pracy (Dz.U. 1974 nr 24 poz. 141; t.j. Dz.U. 2025 poz. 277 ze zm.).
\BibItem{PrAut}{praut} Ustawa z dnia 4 lutego 1994 r. o prawie autorskim i prawach pokrewnych (Dz.U. 1994 nr 24 poz. 83; t.j. Dz.U. 2025 poz. 24 ze zm.).
\BibItem{PWP}{pwp} Ustawa z dnia 30 czerwca 2000 r. --- Prawo własności przemysłowej (Dz.U. 2001 nr 49 poz. 508; t.j. Dz.U. 2023 poz. 1170 ze zm.).
\BibItem{UZNK}{uznk} Ustawa z dnia 16 kwietnia 1993 r. o zwalczaniu nieuczciwej konkurencji (Dz.U. 1993 nr 47 poz. 211; t.j. Dz.U. 2026 poz. 85 ze zm.).
\BibItem{PIT}{pit} Ustawa z dnia 26 lipca 1991 r. o podatku dochodowym od osób fizycznych (Dz.U. 1991 nr 80 poz. 350; t.j. Dz.U. 2026 poz. 592 ze zm.).
\BibItem{CIT}{cit} Ustawa z dnia 15 lutego 1992 r. o podatku dochodowym od osób prawnych (Dz.U. 1992 nr 21 poz. 86 ze zm.).
\BibItem{UoR}{uor} Ustawa z dnia 29 września 1994 r. o rachunkowości (Dz.U. 1994 nr 121 poz. 591; t.j. Dz.U. 2026 poz. 522 ze zm.).
\BibItem{UKRS}{ukrs} Ustawa z dnia 20 sierpnia 1997 r. o Krajowym Rejestrze Sądowym (Dz.U. 1997 nr 121 poz. 769; t.j. Dz.U. 2025 poz. 869 ze zm.).
\BibItem{UONI}{uoni} Ustawa z dnia 5 sierpnia 2010 r. o ochronie informacji niejawnych (Dz.U. 2010 nr 182 poz. 1228; t.j. Dz.U. 2025 poz. 1209 ze zm.).
\BibItem{UOTS}{uots} Ustawa z dnia 29 listopada 2000 r. o obrocie z zagranicą towarami, technologiami i usługami o znaczeniu strategicznym dla bezpieczeństwa państwa, a także dla utrzymania międzynarodowego pokoju i bezpieczeństwa (Dz.U. 2000 nr 119 poz. 1250; t.j. Dz.U. 2023 poz. 1582 ze zm., w tym Dz.U. 2026 poz. 471 --- definicje i zezwolenia, m.in. pomoc techniczna).
\BibItem{UWT}{uwt} Ustawa z dnia 13 czerwca 2019 r. o wykonywaniu działalności gospodarczej w zakresie wytwarzania i obrotu materiałami wybuchowymi, bronią, amunicją oraz wyrobami i technologią o przeznaczeniu wojskowym lub policyjnym (Dz.U. 2019 poz. 1214; t.j. Dz.U. 2023 poz. 1743 ze zm.); koncesja na obrót technologią o przeznaczeniu wojskowym (art.~7), wymogi wobec wspólników z co najmniej 20\% akcji (art.~10 i 17)..
\BibItem{UKI}{uki} Ustawa z dnia 24 lipca 2015 r. o kontroli niektórych inwestycji (Dz.U. 2015 poz. 1272; t.j. Dz.U. 2026 poz. 47 ze zm.).
\BibItem{USank}{usank} Ustawa z dnia 13 kwietnia 2022 r. o szczególnych rozwiązaniach w zakresie przeciwdziałania wspieraniu agresji na Ukrainę oraz służących ochronie bezpieczeństwa narodowego (Dz.U. 2022 poz. 835; t.j. Dz.U. 2025 poz. 514 ze zm.).
\BibItem{USyg}{usyg} Ustawa z dnia 14 czerwca 2024 r. o ochronie sygnalistów (Dz.U. 2024 poz. 928).
\BibItem{UODO}{uodo} Ustawa z dnia 10 maja 2018 r. o ochronie danych osobowych (Dz.U. 2018 poz. 1000 ze zm.).
\BibItem{PrLot}{prlot} Ustawa z dnia 3 lipca 2002 r. --- Prawo lotnicze (Dz.U. 2002 nr 130 poz. 1112; t.j. Dz.U. 2025 poz. 1431 ze zm.).
\BibItem{PKD 2025}{pkd} Rozporządzenie Rady Ministrów z dnia 18 grudnia 2024 r. w sprawie Polskiej Klasyfikacji Działalności (PKD) (Dz.U. 2024 poz. 1936).
\BibItem{WT}{wt} Rozporządzenie Rady Ministrów z dnia 17 września 2019 r. w sprawie klasyfikacji rodzajów materiałów wybuchowych, broni, amunicji oraz wyrobów i technologii o przeznaczeniu wojskowym lub policyjnym, na których wytwarzanie lub obrót jest wymagane uzyskanie koncesji (Dz.U. 2019 poz. 1888) --- obejmuje m.in. bezzałogowe statki powietrzne i ich systemy (WT~V pkt~3).
\end{BibList}

\section*{B. Akty prawne Unii Europejskiej}
\begin{BibList}
\BibItem{RODO}{rodo} Rozporządzenie Parlamentu Europejskiego i Rady (UE) 2016/679 z dnia 27 kwietnia 2016 r. w sprawie ochrony osób fizycznych w związku z przetwarzaniem danych osobowych i w sprawie swobodnego przepływu takich danych oraz uchylenia dyrektywy 95/46/WE (ogólne rozporządzenie o ochronie danych) (Dz.Urz. UE L 119 z 4.5.2016, s.~1).
\BibItem{Dual-use}{dualuse} Rozporządzenie Parlamentu Europejskiego i Rady (UE) 2021/821 z dnia 20 maja 2021 r. ustanawiające unijny system kontroli wywozu, pośrednictwa, pomocy technicznej, tranzytu i transferu produktów podwójnego zastosowania (wersja przekształcona) (Dz.Urz. UE L 206 z 11.6.2021, s.~1, ze zm.).
\BibItem{EDF}{edf} Rozporządzenie Parlamentu Europejskiego i Rady (UE) 2021/697 z dnia 29 kwietnia 2021 r. ustanawiające Europejski Fundusz Obronny i uchylające rozporządzenie (UE) 2018/1092 (Dz.Urz. UE L 170 z 12.5.2021, s.~149).
\BibItem{UAS 945}{uas945} Rozporządzenie delegowane Komisji (UE) 2019/945 z dnia 12 marca 2019 r. w sprawie bezzałogowych systemów powietrznych oraz operatorów bezzałogowych systemów powietrznych z państw trzecich (Dz.Urz. UE L 152 z 11.6.2019, s.~1).
\BibItem{UAS 947}{uas947} Rozporządzenie wykonawcze Komisji (UE) 2019/947 z dnia 24 maja 2019 r. w sprawie przepisów i procedur dotyczących eksploatacji bezzałogowych statków powietrznych (Dz.Urz. UE L 152 z 11.6.2019, s.~45).
\BibItem{FDI}{fdi} Rozporządzenie Parlamentu Europejskiego i Rady (UE) 2019/452 z dnia 19 marca 2019 r. ustanawiające ramy monitorowania bezpośrednich inwestycji zagranicznych w Unii --- uchylone z dniem 17 stycznia 2028 r. przez \Zr{fdi2026}.
\BibItem{FDI 2026}{fdi2026} Rozporządzenie Parlamentu Europejskiego i Rady (UE) 2026/1386 z dnia 17 czerwca 2026 r. w sprawie monitorowania inwestycji zagranicznych w Unii i uchylające rozporządzenie (UE) 2019/452 (Dz.Urz. UE L 2026/1386 z 26.6.2026); stosuje się od 17 stycznia 2028 r. (niektóre przepisy od 16 lipca 2026 r.); obejmuje m.in. spółki rozwijające produkty podwójnego zastosowania, wojskowe i drony.
\BibItem{TTBER}{ttber} Rozporządzenie Komisji (UE) 2026/877 z dnia 16 kwietnia 2026 r. w sprawie stosowania art.~101 ust.~3 TFUE do kategorii porozumień o transferze technologii (Dz.Urz. UE L 2026/877 z 21.4.2026); obowiązuje od 1 maja 2026 r. do 30 kwietnia 2038 r.; zastąpiło rozporządzenie (UE) nr 316/2014 --- licencje na technologię, know-how i oprogramowanie.
\BibItem{AI Act}{aiact} Rozporządzenie Parlamentu Europejskiego i Rady (UE) 2024/1689 z dnia 13 czerwca 2024 r. w sprawie ustanowienia zharmonizowanych przepisów dotyczących sztucznej inteligencji (akt w sprawie sztucznej inteligencji) (Dz.Urz. UE L 2024/1689 z 12.7.2024); systemy wyłącznie wojskowe są wyłączone (art.~2 ust.~3), systemy podwójnego zastosowania --- objęte.
\BibItem{Sankcje}{sankcje} Rozporządzenie Rady (UE) nr 833/2014 z dnia 31 lipca 2014 r. dotyczące środków ograniczających w związku z działaniami Rosji destabilizującymi sytuację na Ukrainie (Dz.Urz. UE L 229 z 31.7.2014, s.~1) oraz rozporządzenie Rady (UE) nr 269/2014 z dnia 17 marca 2014 r. (Dz.Urz. UE L 78 z 17.3.2014, s.~6) (środki ograniczające).
\BibItem{TSD}{tsd} Dyrektywa Parlamentu Europejskiego i Rady (UE) 2016/943 z dnia 8 czerwca 2016 r. w sprawie ochrony niejawnego know-how i niejawnych informacji handlowych (tajemnic przedsiębiorstwa) przed ich bezprawnym pozyskiwaniem, wykorzystywaniem i ujawnianiem (Dz.Urz. UE L 157 z 15.6.2016, s.~1).
\BibItem{WBD}{wbd} Dyrektywa Parlamentu Europejskiego i Rady (UE) 2019/1937 z dnia 23 października 2019 r. w sprawie ochrony osób zgłaszających naruszenia prawa Unii.
\end{BibList}

\section*{C. Orzecznictwo}
\begin{BibList}
\BibItem{SN 61/03}{sn6103} Uchwała składu siedmiu sędziów SN z dnia 6 listopada 2003 r., III CZP 61/03, OSNC 2004, nr 5, poz. 69 (zasada prawna) --- kara umowna należy się także wtedy, gdy wierzyciel nie poniósł szkody; dłużnik może żądać jej miarkowania (art.~484 §~2 KC).
\BibItem{SN 304/00}{sn30400} Wyrok SN z dnia 3 października 2000 r., I CKN 304/00, OSNC 2001, nr 4, poz. 59 --- wykorzystanie przez pracownika informacji, co do których przedsiębiorca nie podjął niezbędnych działań w celu zachowania ich poufności, należy traktować jako wykorzystanie powszechnej wiedzy, do której przedsiębiorca nie ma ustawowych uprawnień.
\BibItem{SN 120/07}{sn12007} Wyrok SN z dnia 8 stycznia 2008 r., II PK 120/07 --- odpowiedzialność pracownika za szkodę regulują wyczerpująco przepisy Kodeksu pracy; kara umowna (art.~483 KC) zastrzeżona w stosunku pracy nie może być dochodzona; dopuszczalna jest natomiast w umowie o zakazie konkurencji po ustaniu stosunku pracy (SN powołuje wyrok z 10 października 2003 r., I PK 528/02, OSNP 2004, nr 19, poz. 336).
\BibItem{SN 42/18}{sn4218} Uchwała SN z dnia 7 września 2018 r., III CZP 42/18, OSNC 2019, nr 6, poz. 64 --- art.~210 §~1 KSH stosuje się do zgody spółki z o.o. na zmianę umowy spółki komandytowej, której wspólnikiem jest także członek zarządu; przepis obejmuje wszystkie czynności związane z umową między spółką a członkiem zarządu (zawarcie, zmianę, wypowiedzenie, rozwiązanie, potrącenie, ugodę), a do jego zastosowania wystarcza potencjalna kolizja interesów. Przydatna przy wykładni art.~300$^{79}$ KSH.
\BibItem{SA 485/18}{sa48518} Wyrok Sądu Apelacyjnego w Warszawie z dnia 9 kwietnia 2019 r., V ACa 485/18 --- członek zarządu nie może jako pełnomocnik reprezentować spółki z o.o. w umowie z innym członkiem zarządu (art.~210 §~1 KSH; sąd powołuje doktrynę oraz uchwały SN III CZP 71/18 i III CZP 42/18); skutkiem naruszenia jest bezwzględna nieważność czynności (art.~58 §~1 KC).
\end{BibList}

\section*{D. Literatura}
\begin{BibList}
\BibItem{Malinowski}{malinowski} P. Malinowski (red.), \emph{Prosta spółka akcyjna. Komentarz}, Warszawa 2021.
\BibItem{Opalski}{opalski} A. Opalski, \emph{Prosta spółka akcyjna --- nowy typ spółki handlowej}, cz.~I: „Przegląd Prawa Handlowego” 2019, nr 11, s.~5--15; cz.~II: tamże, nr 12, s.~4--16.
\BibItem{ECFR}{ecfr} \emph{The Simple Joint Stock Company: Emergence of a New Close Company in Poland}, „European Company and Financial Law Review” 2024, t.~21, DOI 10.1515/ecfr-2024-0021.
\BibItem{Jara}{jara} Z. Jara (red.), \emph{Kodeks spółek handlowych. Komentarz}, C.H. Beck (wydanie aktualne).
\BibItem{Gniewek}{gniewek} E. Gniewek, P. Machnikowski (red.), \emph{Kodeks cywilny. Komentarz}, C.H. Beck (wydanie aktualne).
\BibItem{Barta}{barta} J. Barta, R. Markiewicz (red.), \emph{Ustawa o prawie autorskim i prawach pokrewnych. Komentarz}, Wolters Kluwer (wydanie aktualne).
\BibItem{Kostański}{kostanski} P. Kostański (red.), \emph{Prawo własności przemysłowej. Komentarz}, C.H. Beck (wydanie aktualne).
\BibItem{Szwaja}{szwaja} J. Szwaja (red.), \emph{Ustawa o zwalczaniu nieuczciwej konkurencji. Komentarz}, C.H. Beck (wydanie aktualne).
\end{BibList}

\section*{E. Materiały urzędowe}
\begin{BibList}
\BibItem{GUS PKD}{guspkd} Główny Urząd Statystyczny, \emph{Polska Klasyfikacja Działalności PKD 2025} (klasyfikacje.stat.gov.pl).
\BibItem{GUS CPI}{guscpi} Komunikaty Prezesa GUS w sprawie średniorocznego wskaźnika cen towarów i usług konsumpcyjnych ogółem, ogłaszane w Dzienniku Urzędowym RP „Monitor Polski” (podstawa waloryzacji kar i Cennika).
\end{BibList}

\section*{F. Dokumenty Spółki}
\begin{BibList}
\BibItem{US}{us} Umowa spółki Basilisk Systems P.S.A., wersja 3.4-S (\texttt{psa.pdf}; wersja angielska \texttt{psa\_en.pdf}).
\BibItem{UA}{ua} Umowa akcjonariuszy, wersja 1.8 (\texttt{basilisk\_sha.pdf}).
\BibItem{IP}{ip} Umowy przeniesienia praw autorskich (Umowa nr~1 i nr~2), wersja 1.6 (\texttt{basilisk\_ip.pdf}).
\BibItem{Cennik}{cennik} Cennik wynagrodzeń za rozwiązania projektowe (\texttt{basilisk\_cennik.pdf}).
\BibItem{NDA}{nda} Wzór umowy o zachowaniu poufności (\texttt{basilisk\_nda.pdf}).
\BibItem{Polityka}{polityka} Polityka bezpieczeństwa informacji (\texttt{basilisk\_polityka.pdf}).
\BibItem{Podręcznik}{podrecznik} Podręcznik podejmowania decyzji (\texttt{basilisk\_decyzje.pdf}).
\end{BibList}
 w podręczniku decyzji i komentarzu do umowy spółki.
% Odesłanie w tekście: \Zr{klucz} -> [ETYKIETA] z linkiem do pozycji.
\makeatletter
\providecommand{\BibItem}[2]{\item[\textbf{[#1]}]\def\@currentlabel{#1}\def\@currentlabelname{#1}\phantomsection\label{bib:#2}}
\makeatother
\providecommand{\Zr}[1]{[\ref{bib:#1}]}
\newenvironment{BibList}{\begin{list}{}{\setlength{\labelwidth}{2.6cm}\setlength{\leftmargin}{2.9cm}\setlength{\labelsep}{0.3cm}\setlength{\itemsep}{0.15em}\setlength{\parsep}{0pt}\small}}{\end{list}}

\chapter*{Bibliografia}
\addcontentsline{toc}{chapter}{Bibliografia}

Akty prawne podano według miejsca pierwotnej publikacji oraz ostatniego tekstu jednolitego (stan: listopad 2026), z dopiskiem „ze zm.”; przed użyciem należy sprawdzić, czy nie ogłoszono nowszego tekstu w Internetowym Systemie Aktów Prawnych (isap.sejm.gov.pl) albo w Dzienniku Ustaw (eli.gov.pl). Tezy orzeczeń podano w skrócie; przed powołaniem się na nie w sporze należy sprawdzić pełny tekst.

\section*{A. Akty prawne krajowe}
\begin{BibList}
\BibItem{KSH}{ksh} Ustawa z dnia 15 września 2000 r. --- Kodeks spółek handlowych (Dz.U. 2000 nr 94 poz. 1037 ze zm.; t.j. Dz.U. 2024 poz. 18 ze zm., w tym ustawą z dnia 23 stycznia 2026 r., Dz.U. 2026 poz. 176 --- zmiany rejestru akcjonariuszy od 18 lutego 2027 r.); przepisy o prostej spółce akcyjnej: dział IA, art.~300$^{1}$--300$^{134}$.
\BibItem{Nowela PSA}{psa2019} Ustawa z dnia 19 lipca 2019 r. o zmianie ustawy --- Kodeks spółek handlowych oraz niektórych innych ustaw (Dz.U. 2019 poz. 1655) --- wprowadzenie prostej spółki akcyjnej; ustawa w pierwotnym brzmieniu wchodziła w życie 1 marca 2020 r., a przepisy o PSA stosuje się od 1 lipca 2021 r. (po przesunięciu terminu).
\BibItem{KC}{kc} Ustawa z dnia 23 kwietnia 1964 r. --- Kodeks cywilny (Dz.U. 1964 nr 16 poz. 93; t.j. Dz.U. 2026 poz. 795 ze zm.).
\BibItem{KRO}{kro} Ustawa z dnia 25 lutego 1964 r. --- Kodeks rodzinny i opiekuńczy (Dz.U. 1964 nr 9 poz. 59; t.j. Dz.U. 2026 poz. 236 ze zm.).
\BibItem{KPC}{kpc} Ustawa z dnia 17 listopada 1964 r. --- Kodeks postępowania cywilnego (Dz.U. 1964 nr 43 poz. 296; t.j. Dz.U. 2026 poz. 468 ze zm.).
\BibItem{KP}{kp} Ustawa z dnia 26 czerwca 1974 r. --- Kodeks pracy (Dz.U. 1974 nr 24 poz. 141; t.j. Dz.U. 2025 poz. 277 ze zm.).
\BibItem{PrAut}{praut} Ustawa z dnia 4 lutego 1994 r. o prawie autorskim i prawach pokrewnych (Dz.U. 1994 nr 24 poz. 83; t.j. Dz.U. 2025 poz. 24 ze zm.).
\BibItem{PWP}{pwp} Ustawa z dnia 30 czerwca 2000 r. --- Prawo własności przemysłowej (Dz.U. 2001 nr 49 poz. 508; t.j. Dz.U. 2023 poz. 1170 ze zm.).
\BibItem{UZNK}{uznk} Ustawa z dnia 16 kwietnia 1993 r. o zwalczaniu nieuczciwej konkurencji (Dz.U. 1993 nr 47 poz. 211; t.j. Dz.U. 2026 poz. 85 ze zm.).
\BibItem{PIT}{pit} Ustawa z dnia 26 lipca 1991 r. o podatku dochodowym od osób fizycznych (Dz.U. 1991 nr 80 poz. 350; t.j. Dz.U. 2026 poz. 592 ze zm.).
\BibItem{CIT}{cit} Ustawa z dnia 15 lutego 1992 r. o podatku dochodowym od osób prawnych (Dz.U. 1992 nr 21 poz. 86 ze zm.).
\BibItem{UoR}{uor} Ustawa z dnia 29 września 1994 r. o rachunkowości (Dz.U. 1994 nr 121 poz. 591; t.j. Dz.U. 2026 poz. 522 ze zm.).
\BibItem{UKRS}{ukrs} Ustawa z dnia 20 sierpnia 1997 r. o Krajowym Rejestrze Sądowym (Dz.U. 1997 nr 121 poz. 769; t.j. Dz.U. 2025 poz. 869 ze zm.).
\BibItem{UONI}{uoni} Ustawa z dnia 5 sierpnia 2010 r. o ochronie informacji niejawnych (Dz.U. 2010 nr 182 poz. 1228; t.j. Dz.U. 2025 poz. 1209 ze zm.).
\BibItem{UOTS}{uots} Ustawa z dnia 29 listopada 2000 r. o obrocie z zagranicą towarami, technologiami i usługami o znaczeniu strategicznym dla bezpieczeństwa państwa, a także dla utrzymania międzynarodowego pokoju i bezpieczeństwa (Dz.U. 2000 nr 119 poz. 1250; t.j. Dz.U. 2023 poz. 1582 ze zm., w tym Dz.U. 2026 poz. 471 --- definicje i zezwolenia, m.in. pomoc techniczna).
\BibItem{UWT}{uwt} Ustawa z dnia 13 czerwca 2019 r. o wykonywaniu działalności gospodarczej w zakresie wytwarzania i obrotu materiałami wybuchowymi, bronią, amunicją oraz wyrobami i technologią o przeznaczeniu wojskowym lub policyjnym (Dz.U. 2019 poz. 1214; t.j. Dz.U. 2023 poz. 1743 ze zm.); koncesja na obrót technologią o przeznaczeniu wojskowym (art.~7), wymogi wobec wspólników z co najmniej 20\% akcji (art.~10 i 17)..
\BibItem{UKI}{uki} Ustawa z dnia 24 lipca 2015 r. o kontroli niektórych inwestycji (Dz.U. 2015 poz. 1272; t.j. Dz.U. 2026 poz. 47 ze zm.).
\BibItem{USank}{usank} Ustawa z dnia 13 kwietnia 2022 r. o szczególnych rozwiązaniach w zakresie przeciwdziałania wspieraniu agresji na Ukrainę oraz służących ochronie bezpieczeństwa narodowego (Dz.U. 2022 poz. 835; t.j. Dz.U. 2025 poz. 514 ze zm.).
\BibItem{USyg}{usyg} Ustawa z dnia 14 czerwca 2024 r. o ochronie sygnalistów (Dz.U. 2024 poz. 928).
\BibItem{UODO}{uodo} Ustawa z dnia 10 maja 2018 r. o ochronie danych osobowych (Dz.U. 2018 poz. 1000 ze zm.).
\BibItem{PrLot}{prlot} Ustawa z dnia 3 lipca 2002 r. --- Prawo lotnicze (Dz.U. 2002 nr 130 poz. 1112; t.j. Dz.U. 2025 poz. 1431 ze zm.).
\BibItem{PKD 2025}{pkd} Rozporządzenie Rady Ministrów z dnia 18 grudnia 2024 r. w sprawie Polskiej Klasyfikacji Działalności (PKD) (Dz.U. 2024 poz. 1936).
\BibItem{WT}{wt} Rozporządzenie Rady Ministrów z dnia 17 września 2019 r. w sprawie klasyfikacji rodzajów materiałów wybuchowych, broni, amunicji oraz wyrobów i technologii o przeznaczeniu wojskowym lub policyjnym, na których wytwarzanie lub obrót jest wymagane uzyskanie koncesji (Dz.U. 2019 poz. 1888) --- obejmuje m.in. bezzałogowe statki powietrzne i ich systemy (WT~V pkt~3).
\end{BibList}

\section*{B. Akty prawne Unii Europejskiej}
\begin{BibList}
\BibItem{RODO}{rodo} Rozporządzenie Parlamentu Europejskiego i Rady (UE) 2016/679 z dnia 27 kwietnia 2016 r. w sprawie ochrony osób fizycznych w związku z przetwarzaniem danych osobowych i w sprawie swobodnego przepływu takich danych oraz uchylenia dyrektywy 95/46/WE (ogólne rozporządzenie o ochronie danych) (Dz.Urz. UE L 119 z 4.5.2016, s.~1).
\BibItem{Dual-use}{dualuse} Rozporządzenie Parlamentu Europejskiego i Rady (UE) 2021/821 z dnia 20 maja 2021 r. ustanawiające unijny system kontroli wywozu, pośrednictwa, pomocy technicznej, tranzytu i transferu produktów podwójnego zastosowania (wersja przekształcona) (Dz.Urz. UE L 206 z 11.6.2021, s.~1, ze zm.).
\BibItem{EDF}{edf} Rozporządzenie Parlamentu Europejskiego i Rady (UE) 2021/697 z dnia 29 kwietnia 2021 r. ustanawiające Europejski Fundusz Obronny i uchylające rozporządzenie (UE) 2018/1092 (Dz.Urz. UE L 170 z 12.5.2021, s.~149).
\BibItem{UAS 945}{uas945} Rozporządzenie delegowane Komisji (UE) 2019/945 z dnia 12 marca 2019 r. w sprawie bezzałogowych systemów powietrznych oraz operatorów bezzałogowych systemów powietrznych z państw trzecich (Dz.Urz. UE L 152 z 11.6.2019, s.~1).
\BibItem{UAS 947}{uas947} Rozporządzenie wykonawcze Komisji (UE) 2019/947 z dnia 24 maja 2019 r. w sprawie przepisów i procedur dotyczących eksploatacji bezzałogowych statków powietrznych (Dz.Urz. UE L 152 z 11.6.2019, s.~45).
\BibItem{FDI}{fdi} Rozporządzenie Parlamentu Europejskiego i Rady (UE) 2019/452 z dnia 19 marca 2019 r. ustanawiające ramy monitorowania bezpośrednich inwestycji zagranicznych w Unii --- uchylone z dniem 17 stycznia 2028 r. przez \Zr{fdi2026}.
\BibItem{FDI 2026}{fdi2026} Rozporządzenie Parlamentu Europejskiego i Rady (UE) 2026/1386 z dnia 17 czerwca 2026 r. w sprawie monitorowania inwestycji zagranicznych w Unii i uchylające rozporządzenie (UE) 2019/452 (Dz.Urz. UE L 2026/1386 z 26.6.2026); stosuje się od 17 stycznia 2028 r. (niektóre przepisy od 16 lipca 2026 r.); obejmuje m.in. spółki rozwijające produkty podwójnego zastosowania, wojskowe i drony.
\BibItem{TTBER}{ttber} Rozporządzenie Komisji (UE) 2026/877 z dnia 16 kwietnia 2026 r. w sprawie stosowania art.~101 ust.~3 TFUE do kategorii porozumień o transferze technologii (Dz.Urz. UE L 2026/877 z 21.4.2026); obowiązuje od 1 maja 2026 r. do 30 kwietnia 2038 r.; zastąpiło rozporządzenie (UE) nr 316/2014 --- licencje na technologię, know-how i oprogramowanie.
\BibItem{AI Act}{aiact} Rozporządzenie Parlamentu Europejskiego i Rady (UE) 2024/1689 z dnia 13 czerwca 2024 r. w sprawie ustanowienia zharmonizowanych przepisów dotyczących sztucznej inteligencji (akt w sprawie sztucznej inteligencji) (Dz.Urz. UE L 2024/1689 z 12.7.2024); systemy wyłącznie wojskowe są wyłączone (art.~2 ust.~3), systemy podwójnego zastosowania --- objęte.
\BibItem{Sankcje}{sankcje} Rozporządzenie Rady (UE) nr 833/2014 z dnia 31 lipca 2014 r. dotyczące środków ograniczających w związku z działaniami Rosji destabilizującymi sytuację na Ukrainie (Dz.Urz. UE L 229 z 31.7.2014, s.~1) oraz rozporządzenie Rady (UE) nr 269/2014 z dnia 17 marca 2014 r. (Dz.Urz. UE L 78 z 17.3.2014, s.~6) (środki ograniczające).
\BibItem{TSD}{tsd} Dyrektywa Parlamentu Europejskiego i Rady (UE) 2016/943 z dnia 8 czerwca 2016 r. w sprawie ochrony niejawnego know-how i niejawnych informacji handlowych (tajemnic przedsiębiorstwa) przed ich bezprawnym pozyskiwaniem, wykorzystywaniem i ujawnianiem (Dz.Urz. UE L 157 z 15.6.2016, s.~1).
\BibItem{WBD}{wbd} Dyrektywa Parlamentu Europejskiego i Rady (UE) 2019/1937 z dnia 23 października 2019 r. w sprawie ochrony osób zgłaszających naruszenia prawa Unii.
\end{BibList}

\section*{C. Orzecznictwo}
\begin{BibList}
\BibItem{SN 61/03}{sn6103} Uchwała składu siedmiu sędziów SN z dnia 6 listopada 2003 r., III CZP 61/03, OSNC 2004, nr 5, poz. 69 (zasada prawna) --- kara umowna należy się także wtedy, gdy wierzyciel nie poniósł szkody; dłużnik może żądać jej miarkowania (art.~484 §~2 KC).
\BibItem{SN 304/00}{sn30400} Wyrok SN z dnia 3 października 2000 r., I CKN 304/00, OSNC 2001, nr 4, poz. 59 --- wykorzystanie przez pracownika informacji, co do których przedsiębiorca nie podjął niezbędnych działań w celu zachowania ich poufności, należy traktować jako wykorzystanie powszechnej wiedzy, do której przedsiębiorca nie ma ustawowych uprawnień.
\BibItem{SN 120/07}{sn12007} Wyrok SN z dnia 8 stycznia 2008 r., II PK 120/07 --- odpowiedzialność pracownika za szkodę regulują wyczerpująco przepisy Kodeksu pracy; kara umowna (art.~483 KC) zastrzeżona w stosunku pracy nie może być dochodzona; dopuszczalna jest natomiast w umowie o zakazie konkurencji po ustaniu stosunku pracy (SN powołuje wyrok z 10 października 2003 r., I PK 528/02, OSNP 2004, nr 19, poz. 336).
\BibItem{SN 42/18}{sn4218} Uchwała SN z dnia 7 września 2018 r., III CZP 42/18, OSNC 2019, nr 6, poz. 64 --- art.~210 §~1 KSH stosuje się do zgody spółki z o.o. na zmianę umowy spółki komandytowej, której wspólnikiem jest także członek zarządu; przepis obejmuje wszystkie czynności związane z umową między spółką a członkiem zarządu (zawarcie, zmianę, wypowiedzenie, rozwiązanie, potrącenie, ugodę), a do jego zastosowania wystarcza potencjalna kolizja interesów. Przydatna przy wykładni art.~300$^{79}$ KSH.
\BibItem{SA 485/18}{sa48518} Wyrok Sądu Apelacyjnego w Warszawie z dnia 9 kwietnia 2019 r., V ACa 485/18 --- członek zarządu nie może jako pełnomocnik reprezentować spółki z o.o. w umowie z innym członkiem zarządu (art.~210 §~1 KSH; sąd powołuje doktrynę oraz uchwały SN III CZP 71/18 i III CZP 42/18); skutkiem naruszenia jest bezwzględna nieważność czynności (art.~58 §~1 KC).
\end{BibList}

\section*{D. Literatura}
\begin{BibList}
\BibItem{Malinowski}{malinowski} P. Malinowski (red.), \emph{Prosta spółka akcyjna. Komentarz}, Warszawa 2021.
\BibItem{Opalski}{opalski} A. Opalski, \emph{Prosta spółka akcyjna --- nowy typ spółki handlowej}, cz.~I: „Przegląd Prawa Handlowego” 2019, nr 11, s.~5--15; cz.~II: tamże, nr 12, s.~4--16.
\BibItem{ECFR}{ecfr} \emph{The Simple Joint Stock Company: Emergence of a New Close Company in Poland}, „European Company and Financial Law Review” 2024, t.~21, DOI 10.1515/ecfr-2024-0021.
\BibItem{Jara}{jara} Z. Jara (red.), \emph{Kodeks spółek handlowych. Komentarz}, C.H. Beck (wydanie aktualne).
\BibItem{Gniewek}{gniewek} E. Gniewek, P. Machnikowski (red.), \emph{Kodeks cywilny. Komentarz}, C.H. Beck (wydanie aktualne).
\BibItem{Barta}{barta} J. Barta, R. Markiewicz (red.), \emph{Ustawa o prawie autorskim i prawach pokrewnych. Komentarz}, Wolters Kluwer (wydanie aktualne).
\BibItem{Kostański}{kostanski} P. Kostański (red.), \emph{Prawo własności przemysłowej. Komentarz}, C.H. Beck (wydanie aktualne).
\BibItem{Szwaja}{szwaja} J. Szwaja (red.), \emph{Ustawa o zwalczaniu nieuczciwej konkurencji. Komentarz}, C.H. Beck (wydanie aktualne).
\end{BibList}

\section*{E. Materiały urzędowe}
\begin{BibList}
\BibItem{GUS PKD}{guspkd} Główny Urząd Statystyczny, \emph{Polska Klasyfikacja Działalności PKD 2025} (klasyfikacje.stat.gov.pl).
\BibItem{GUS CPI}{guscpi} Komunikaty Prezesa GUS w sprawie średniorocznego wskaźnika cen towarów i usług konsumpcyjnych ogółem, ogłaszane w Dzienniku Urzędowym RP „Monitor Polski” (podstawa waloryzacji kar i Cennika).
\end{BibList}

\section*{F. Dokumenty Spółki}
\begin{BibList}
\BibItem{US}{us} Umowa spółki Basilisk Systems P.S.A., wersja 3.4-S (\texttt{psa.pdf}; wersja angielska \texttt{psa\_en.pdf}).
\BibItem{UA}{ua} Umowa akcjonariuszy, wersja 1.8 (\texttt{basilisk\_sha.pdf}).
\BibItem{IP}{ip} Umowy przeniesienia praw autorskich (Umowa nr~1 i nr~2), wersja 1.6 (\texttt{basilisk\_ip.pdf}).
\BibItem{Cennik}{cennik} Cennik wynagrodzeń za rozwiązania projektowe (\texttt{basilisk\_cennik.pdf}).
\BibItem{NDA}{nda} Wzór umowy o zachowaniu poufności (\texttt{basilisk\_nda.pdf}).
\BibItem{Polityka}{polityka} Polityka bezpieczeństwa informacji (\texttt{basilisk\_polityka.pdf}).
\BibItem{Podręcznik}{podrecznik} Podręcznik podejmowania decyzji (\texttt{basilisk\_decyzje.pdf}).
\end{BibList}


\end{document}
